\chapter{State-Space Representation}
\label{ch:state_space_representation}

\whythismatters{
Transfer functions have served us well for analyzing and designing single-input single-output (SISO) control systems. However, as you advance in control engineering, you will encounter systems that transfer functions cannot adequately describe: aircraft with multiple control surfaces affecting multiple outputs, robotic manipulators with coupled joint dynamics, chemical processes with interacting temperature and concentration loops. These multi-input multi-output (MIMO) systems demand a more powerful mathematical framework.

State-space representation provides that framework. Instead of describing only the input-output relationship, state-space models capture the complete internal dynamics of a system through a set of first-order differential equations. This approach reveals properties that transfer functions hide---whether all system modes can be controlled, whether internal states can be estimated from measurements, and how energy flows through the system.

More importantly, state-space representation is the foundation for virtually all modern control techniques: optimal control (LQR), state estimation (Kalman filtering), model predictive control (MPC), and the machine learning methods we will explore in later chapters. Mastering state-space methods transforms you from a classical control engineer limited to SISO systems into a modern control engineer capable of tackling the complex multivariable systems that define contemporary engineering challenges.

By the end of this chapter, you will derive state-space models from physical systems, convert between transfer functions and state-space forms, analyze controllability and observability, transform between canonical forms, and compute state transition matrices---the essential skills for modern control design.
}

%=============================================================================
\section{State-Space Fundamentals}
\label{sec:state_space_fundamentals}
%=============================================================================

I want to begin our exploration of state-space methods by establishing the fundamental concepts that distinguish this approach from the transfer function methods we have studied previously. The key insight is that state-space representation captures \textit{internal} system dynamics, not just input-output behavior.

\subsection{State Variables and the State Vector}

\begin{definition}[State Variables]
\textbf{State variables} are the smallest set of variables that completely characterize the dynamic behavior of a system at any time $t$. Given the state variables at time $t_0$ and the input for $t \geq t_0$, we can determine the state and output for all future time.
\end{definition}

The number of state variables required equals the number of independent energy storage elements in the system. This is a profound connection between physics and mathematics:

\begin{center}
\begin{tabular}{lll}
\toprule
\textbf{System Type} & \textbf{Energy Storage Element} & \textbf{State Variable} \\
\midrule
Mechanical (translational) & Mass & Position, velocity \\
Mechanical (rotational) & Inertia & Angle, angular velocity \\
Electrical & Capacitor & Voltage across capacitor \\
Electrical & Inductor & Current through inductor \\
Thermal & Thermal mass & Temperature \\
Hydraulic & Fluid capacitance & Pressure or volume \\
\bottomrule
\end{tabular}
\end{center}

We collect all state variables into the \textbf{state vector}:
\begin{equation}
\mathbf{x}(t) = \begin{bmatrix} x_1(t) \\ x_2(t) \\ \vdots \\ x_n(t) \end{bmatrix}
\end{equation}
where $n$ is the \textbf{order} of the system. The state vector lives in an $n$-dimensional \textbf{state space}, and as time evolves, it traces a trajectory through this space.

\begin{example}[State Variables for a Mass-Spring-Damper System]
\label{ex:msd_state_variables}
Consider a mass $m$ connected to a spring (stiffness $k$) and damper (coefficient $b$), with applied force $F(t)$. Newton's second law gives:
\begin{equation}
m\ddot{x} + b\dot{x} + kx = F
\end{equation}

This second-order equation has two energy storage elements:
\begin{itemize}
    \item The spring stores potential energy: $E_{\text{spring}} = \frac{1}{2}kx^2$
    \item The mass stores kinetic energy: $E_{\text{kinetic}} = \frac{1}{2}m\dot{x}^2$
\end{itemize}
The damper dissipates energy but does not store it.

Therefore, we need two state variables. The natural choices are:
\begin{equation}
x_1 = x \quad \text{(position)}, \qquad x_2 = \dot{x} \quad \text{(velocity)}
\end{equation}

The state vector is:
\begin{equation}
\mathbf{x} = \begin{bmatrix} x_1 \\ x_2 \end{bmatrix} = \begin{bmatrix} \text{position} \\ \text{velocity} \end{bmatrix}
\end{equation}

Given position and velocity at any instant, plus the future force input, we can completely determine all future motion.
\end{example}

\subsection{State Equations}

The \textbf{state equation} describes how the state vector evolves over time. For a general nonlinear system:
\begin{equation}
\dot{\mathbf{x}}(t) = \mathbf{f}(\mathbf{x}(t), \mathbf{u}(t), t)
\end{equation}
where $\mathbf{u}(t)$ is the input vector.

For \textbf{linear time-invariant (LTI)} systems, this simplifies to:
\begin{equation}
\boxed{\dot{\mathbf{x}}(t) = \mathbf{A}\mathbf{x}(t) + \mathbf{B}\mathbf{u}(t)}
\label{eq:state_equation}
\end{equation}

This is the \textbf{state equation}. Let me explain each term:
\begin{itemize}
    \item $\dot{\mathbf{x}}(t)$: The derivative of the state vector---how the state changes over time
    \item $\mathbf{A}\mathbf{x}(t)$: Contribution from the current state (internal dynamics)
    \item $\mathbf{B}\mathbf{u}(t)$: Contribution from the input (external forcing)
\end{itemize}

\begin{example}[State Equations for Mass-Spring-Damper]
Continuing Example~\ref{ex:msd_state_variables}, let us derive the state equations.

From the definition $x_1 = x$ and $x_2 = \dot{x}$:
\begin{equation}
\dot{x}_1 = \dot{x} = x_2
\end{equation}

From Newton's law $m\ddot{x} = F - b\dot{x} - kx$:
\begin{equation}
\dot{x}_2 = \ddot{x} = -\frac{k}{m}x - \frac{b}{m}\dot{x} + \frac{1}{m}F = -\frac{k}{m}x_1 - \frac{b}{m}x_2 + \frac{1}{m}F
\end{equation}

In matrix form:
\begin{equation}
\begin{bmatrix} \dot{x}_1 \\ \dot{x}_2 \end{bmatrix} = \begin{bmatrix} 0 & 1 \\ -\frac{k}{m} & -\frac{b}{m} \end{bmatrix} \begin{bmatrix} x_1 \\ x_2 \end{bmatrix} + \begin{bmatrix} 0 \\ \frac{1}{m} \end{bmatrix} F
\end{equation}

Thus:
\begin{equation}
\mathbf{A} = \begin{bmatrix} 0 & 1 \\ -\frac{k}{m} & -\frac{b}{m} \end{bmatrix}, \qquad \mathbf{B} = \begin{bmatrix} 0 \\ \frac{1}{m} \end{bmatrix}
\end{equation}
\end{example}

\subsection{Output Equations}

Not all state variables are directly measurable. The \textbf{output equation} relates the measured outputs to the states and inputs:
\begin{equation}
\boxed{\mathbf{y}(t) = \mathbf{C}\mathbf{x}(t) + \mathbf{D}\mathbf{u}(t)}
\label{eq:output_equation}
\end{equation}

where:
\begin{itemize}
    \item $\mathbf{y}(t) \in \mathbb{R}^p$ is the output vector ($p$ outputs)
    \item $\mathbf{C} \in \mathbb{R}^{p \times n}$ is the output matrix
    \item $\mathbf{D} \in \mathbb{R}^{p \times m}$ is the feedthrough (or direct transmission) matrix
\end{itemize}

\textbf{Strictly proper systems} have $\mathbf{D} = \mathbf{0}$, meaning there is no instantaneous transmission from input to output. Most physical systems are strictly proper because energy storage elements prevent instantaneous response.

\subsection{Complete State-Space Representation}

Combining equations~\eqref{eq:state_equation} and~\eqref{eq:output_equation}, the complete state-space representation is:
\begin{empheq}[box=\fbox]{align}
\dot{\mathbf{x}}(t) &= \mathbf{A}\mathbf{x}(t) + \mathbf{B}\mathbf{u}(t) \label{eq:ss_state}\\
\mathbf{y}(t) &= \mathbf{C}\mathbf{x}(t) + \mathbf{D}\mathbf{u}(t) \label{eq:ss_output}
\end{empheq}

The four matrices $(\mathbf{A}, \mathbf{B}, \mathbf{C}, \mathbf{D})$ completely characterize the LTI system:

\begin{center}
\begin{tabular}{ccll}
\toprule
\textbf{Matrix} & \textbf{Dimension} & \textbf{Name} & \textbf{Physical Meaning} \\
\midrule
$\mathbf{A}$ & $n \times n$ & State matrix & Internal dynamics, coupling between states \\
$\mathbf{B}$ & $n \times m$ & Input matrix & How inputs affect state derivatives \\
$\mathbf{C}$ & $p \times n$ & Output matrix & Which states are measured \\
$\mathbf{D}$ & $p \times m$ & Feedthrough matrix & Direct input-to-output transmission \\
\bottomrule
\end{tabular}
\end{center}

where $n$ = number of states, $m$ = number of inputs, $p$ = number of outputs.

\subsection{Physical Interpretation of the System Matrices}

Understanding the physical meaning of each matrix element deepens your intuition for state-space models.

\subsubsection{The A Matrix: Internal Dynamics}

The element $a_{ij}$ of the $\mathbf{A}$ matrix represents the influence of state $x_j$ on the rate of change of state $x_i$:
\begin{equation}
\dot{x}_i = a_{i1}x_1 + a_{i2}x_2 + \cdots + a_{in}x_n + \text{(input terms)}
\end{equation}

\begin{itemize}
    \item \textbf{Diagonal elements} $a_{ii}$: Self-dynamics---how a state affects its own derivative. Negative values indicate damping or decay.
    \item \textbf{Off-diagonal elements} $a_{ij}$ ($i \neq j$): Coupling---how state $x_j$ influences the evolution of state $x_i$.
\end{itemize}

The \textbf{eigenvalues} of $\mathbf{A}$ are the system poles, determining stability and natural response modes.

\subsubsection{The B Matrix: Control Influence}

The element $b_{ij}$ indicates how input $u_j$ affects the derivative of state $x_i$:
\begin{equation}
\dot{x}_i = \text{(state terms)} + b_{i1}u_1 + b_{i2}u_2 + \cdots + b_{im}u_m
\end{equation}

Column $j$ of $\mathbf{B}$ shows which states are directly affected by input $u_j$.

\subsubsection{The C Matrix: Measurement Configuration}

Row $i$ of $\mathbf{C}$ defines how the states combine to form output $y_i$:
\begin{equation}
y_i = c_{i1}x_1 + c_{i2}x_2 + \cdots + c_{in}x_n + \text{(feedthrough terms)}
\end{equation}

Common cases:
\begin{itemize}
    \item $\mathbf{C} = \begin{bmatrix} 1 & 0 \end{bmatrix}$: Only the first state is measured
    \item $\mathbf{C} = \mathbf{I}$: All states are directly measured (full state feedback possible)
\end{itemize}

\subsection{SISO vs. MIMO Systems}

The beauty of state-space representation is that it handles SISO and MIMO systems with the same mathematical framework. The only difference is matrix dimensions:

\begin{center}
\begin{tabular}{lcccc}
\toprule
\textbf{System Type} & $\mathbf{A}$ & $\mathbf{B}$ & $\mathbf{C}$ & $\mathbf{D}$ \\
\midrule
SISO ($m=1$, $p=1$) & $n \times n$ & $n \times 1$ (vector) & $1 \times n$ (vector) & scalar \\
MIMO ($m$ inputs, $p$ outputs) & $n \times n$ & $n \times m$ & $p \times n$ & $p \times m$ \\
\bottomrule
\end{tabular}
\end{center}

This unified treatment is a major advantage over transfer functions, where MIMO systems require cumbersome transfer function matrices.

\subsection{Connection to Transfer Functions}

For SISO systems, the transfer function can be derived from the state-space representation. Taking Laplace transforms of equations~\eqref{eq:ss_state} and~\eqref{eq:ss_output} with zero initial conditions:
\begin{align}
s\mathbf{X}(s) &= \mathbf{A}\mathbf{X}(s) + \mathbf{B}U(s) \\
Y(s) &= \mathbf{C}\mathbf{X}(s) + \mathbf{D}U(s)
\end{align}

Solving for $\mathbf{X}(s)$:
\begin{equation}
(s\mathbf{I} - \mathbf{A})\mathbf{X}(s) = \mathbf{B}U(s) \implies \mathbf{X}(s) = (s\mathbf{I} - \mathbf{A})^{-1}\mathbf{B}U(s)
\end{equation}

Substituting into the output equation:
\begin{equation}
Y(s) = \mathbf{C}(s\mathbf{I} - \mathbf{A})^{-1}\mathbf{B}U(s) + \mathbf{D}U(s)
\end{equation}

Therefore, the transfer function is:
\begin{equation}
\boxed{G(s) = \frac{Y(s)}{U(s)} = \mathbf{C}(s\mathbf{I} - \mathbf{A})^{-1}\mathbf{B} + \mathbf{D}}
\label{eq:ss_to_tf}
\end{equation}

This fundamental relationship bridges state-space and transfer function representations.

\begin{theorem}[Poles and Eigenvalues]
The poles of the transfer function $G(s)$ are the eigenvalues of the state matrix $\mathbf{A}$.
\end{theorem}

\begin{proof}
The transfer function is $G(s) = \mathbf{C}(s\mathbf{I} - \mathbf{A})^{-1}\mathbf{B} + \mathbf{D}$. Using the formula for matrix inverse:
\begin{equation}
(s\mathbf{I} - \mathbf{A})^{-1} = \frac{\text{adj}(s\mathbf{I} - \mathbf{A})}{\det(s\mathbf{I} - \mathbf{A})}
\end{equation}
The denominator is $\det(s\mathbf{I} - \mathbf{A})$, which is the characteristic polynomial of $\mathbf{A}$. Its roots are the eigenvalues of $\mathbf{A}$, which appear as poles in $G(s)$ (unless cancelled by numerator zeros).
\end{proof}

%=============================================================================
\section{Converting Transfer Functions to State-Space}
\label{sec:tf_to_ss}
%=============================================================================

Given a transfer function, we often need to convert it to state-space form for modern control design. This conversion is not unique---infinitely many state-space realizations exist for any transfer function. However, several \textbf{canonical forms} provide systematic, interpretable representations.

\subsection{Controllable Canonical Form}

The \textbf{controllable canonical form} (also called \textit{phase-variable} or \textit{companion} form) is the most commonly used realization for control design.

Consider a strictly proper transfer function:
\begin{equation}
G(s) = \frac{b_{n-1}s^{n-1} + b_{n-2}s^{n-2} + \cdots + b_1s + b_0}{s^n + a_{n-1}s^{n-1} + \cdots + a_1s + a_0}
\label{eq:tf_general}
\end{equation}

The \textbf{controllable canonical form} realization is:
\begin{equation}
\mathbf{A}_c = \begin{bmatrix}
0 & 1 & 0 & \cdots & 0 \\
0 & 0 & 1 & \cdots & 0 \\
\vdots & \vdots & \vdots & \ddots & \vdots \\
0 & 0 & 0 & \cdots & 1 \\
-a_0 & -a_1 & -a_2 & \cdots & -a_{n-1}
\end{bmatrix}, \quad
\mathbf{B}_c = \begin{bmatrix} 0 \\ 0 \\ \vdots \\ 0 \\ 1 \end{bmatrix}
\label{eq:ccf_AB}
\end{equation}

\begin{equation}
\mathbf{C}_c = \begin{bmatrix} b_0 & b_1 & b_2 & \cdots & b_{n-1} \end{bmatrix}, \quad \mathbf{D}_c = 0
\label{eq:ccf_CD}
\end{equation}

\textbf{Key properties of controllable canonical form:}
\begin{enumerate}
    \item The characteristic polynomial coefficients appear in the \textbf{last row} of $\mathbf{A}_c$ (negated).
    \item The $\mathbf{B}_c$ matrix has a simple structure: zeros with a 1 in the last position.
    \item The system is \textbf{always controllable} (the controllability matrix has full rank).
    \item The $\mathbf{C}_c$ matrix contains the numerator polynomial coefficients.
\end{enumerate}

\begin{example}[Second-Order Transfer Function to CCF]
\label{ex:tf_to_ccf_2nd}
Convert the following transfer function to controllable canonical form:
\begin{equation}
G(s) = \frac{2s + 3}{s^2 + 5s + 6}
\end{equation}

\textbf{Step 1: Identify coefficients}

Comparing with equation~\eqref{eq:tf_general} for $n = 2$:
\begin{itemize}
    \item Denominator: $s^2 + 5s + 6 \implies a_1 = 5, \, a_0 = 6$
    \item Numerator: $2s + 3 \implies b_1 = 2, \, b_0 = 3$
\end{itemize}

\textbf{Step 2: Construct the matrices}

Using equations~\eqref{eq:ccf_AB} and~\eqref{eq:ccf_CD}:
\begin{equation}
\mathbf{A}_c = \begin{bmatrix} 0 & 1 \\ -6 & -5 \end{bmatrix}, \quad
\mathbf{B}_c = \begin{bmatrix} 0 \\ 1 \end{bmatrix}, \quad
\mathbf{C}_c = \begin{bmatrix} 3 & 2 \end{bmatrix}, \quad
\mathbf{D}_c = 0
\end{equation}

\textbf{Step 3: Verify the transfer function}

Let us confirm this realization produces the correct transfer function using equation~\eqref{eq:ss_to_tf}.

First, compute $s\mathbf{I} - \mathbf{A}_c$:
\begin{equation}
s\mathbf{I} - \mathbf{A}_c = \begin{bmatrix} s & -1 \\ 6 & s+5 \end{bmatrix}
\end{equation}

The determinant (characteristic polynomial):
\begin{equation}
\det(s\mathbf{I} - \mathbf{A}_c) = s(s+5) + 6 = s^2 + 5s + 6 \quad \checkmark
\end{equation}

The adjugate matrix:
\begin{equation}
\text{adj}(s\mathbf{I} - \mathbf{A}_c) = \begin{bmatrix} s+5 & 1 \\ -6 & s \end{bmatrix}
\end{equation}

The inverse:
\begin{equation}
(s\mathbf{I} - \mathbf{A}_c)^{-1} = \frac{1}{s^2 + 5s + 6}\begin{bmatrix} s+5 & 1 \\ -6 & s \end{bmatrix}
\end{equation}

Now compute the transfer function:
\begin{align}
G(s) &= \mathbf{C}_c(s\mathbf{I} - \mathbf{A}_c)^{-1}\mathbf{B}_c \\
&= \begin{bmatrix} 3 & 2 \end{bmatrix} \cdot \frac{1}{s^2 + 5s + 6}\begin{bmatrix} s+5 & 1 \\ -6 & s \end{bmatrix} \begin{bmatrix} 0 \\ 1 \end{bmatrix} \\
&= \begin{bmatrix} 3 & 2 \end{bmatrix} \cdot \frac{1}{s^2 + 5s + 6}\begin{bmatrix} 1 \\ s \end{bmatrix} \\
&= \frac{3 + 2s}{s^2 + 5s + 6} = \frac{2s + 3}{s^2 + 5s + 6} \quad \checkmark
\end{align}
\end{example}

\begin{example}[Third-Order Transfer Function to CCF]
\label{ex:tf_to_ccf_3rd}
Convert to controllable canonical form:
\begin{equation}
G(s) = \frac{4s^2 + 3s + 2}{s^3 + 6s^2 + 11s + 6}
\end{equation}

\textbf{Step 1: Identify coefficients}

For $n = 3$:
\begin{itemize}
    \item Denominator: $a_2 = 6, \, a_1 = 11, \, a_0 = 6$
    \item Numerator: $b_2 = 4, \, b_1 = 3, \, b_0 = 2$
\end{itemize}

\textbf{Step 2: Construct the matrices}
\begin{equation}
\mathbf{A}_c = \begin{bmatrix} 0 & 1 & 0 \\ 0 & 0 & 1 \\ -6 & -11 & -6 \end{bmatrix}, \quad
\mathbf{B}_c = \begin{bmatrix} 0 \\ 0 \\ 1 \end{bmatrix}
\end{equation}
\begin{equation}
\mathbf{C}_c = \begin{bmatrix} 2 & 3 & 4 \end{bmatrix}, \quad \mathbf{D}_c = 0
\end{equation}

\textbf{Step 3: Verify poles}

The characteristic polynomial is:
\begin{equation}
\det(s\mathbf{I} - \mathbf{A}_c) = s^3 + 6s^2 + 11s + 6 = (s+1)(s+2)(s+3)
\end{equation}

The poles are $s = -1, -2, -3$, which are the eigenvalues of $\mathbf{A}_c$.
\end{example}

\subsection{Observable Canonical Form}

The \textbf{observable canonical form} is the dual of the controllable form, designed to guarantee observability.

For the same transfer function~\eqref{eq:tf_general}:
\begin{equation}
\mathbf{A}_o = \begin{bmatrix}
0 & 0 & \cdots & 0 & -a_0 \\
1 & 0 & \cdots & 0 & -a_1 \\
0 & 1 & \cdots & 0 & -a_2 \\
\vdots & \vdots & \ddots & \vdots & \vdots \\
0 & 0 & \cdots & 1 & -a_{n-1}
\end{bmatrix}, \quad
\mathbf{B}_o = \begin{bmatrix} b_0 \\ b_1 \\ b_2 \\ \vdots \\ b_{n-1} \end{bmatrix}
\label{eq:ocf_AB}
\end{equation}

\begin{equation}
\mathbf{C}_o = \begin{bmatrix} 0 & 0 & \cdots & 0 & 1 \end{bmatrix}, \quad \mathbf{D}_o = 0
\label{eq:ocf_CD}
\end{equation}

\textbf{Key properties of observable canonical form:}
\begin{enumerate}
    \item The characteristic polynomial coefficients appear in the \textbf{last column} of $\mathbf{A}_o$ (negated).
    \item The $\mathbf{C}_o$ matrix has a simple structure: zeros with a 1 in the last position.
    \item The system is \textbf{always observable} (the observability matrix has full rank).
    \item Observable canonical form is the \textbf{transpose dual} of controllable canonical form:
    \begin{equation}
    \mathbf{A}_o = \mathbf{A}_c^\top, \quad \mathbf{B}_o = \mathbf{C}_c^\top, \quad \mathbf{C}_o = \mathbf{B}_c^\top
    \end{equation}
\end{enumerate}

\begin{example}[Second-Order Transfer Function to OCF]
For the same transfer function from Example~\ref{ex:tf_to_ccf_2nd}:
\begin{equation}
G(s) = \frac{2s + 3}{s^2 + 5s + 6}
\end{equation}

The observable canonical form is:
\begin{equation}
\mathbf{A}_o = \begin{bmatrix} 0 & -6 \\ 1 & -5 \end{bmatrix}, \quad
\mathbf{B}_o = \begin{bmatrix} 3 \\ 2 \end{bmatrix}, \quad
\mathbf{C}_o = \begin{bmatrix} 0 & 1 \end{bmatrix}, \quad
\mathbf{D}_o = 0
\end{equation}

Note that $\mathbf{A}_o = \mathbf{A}_c^\top$, $\mathbf{B}_o = \mathbf{C}_c^\top$, and $\mathbf{C}_o = \mathbf{B}_c^\top$.
\end{example}

\subsection{Diagonal (Modal) Canonical Form}

When the transfer function has \textbf{distinct poles}, we can obtain a \textbf{diagonal canonical form} where the system dynamics are completely decoupled.

\textbf{Procedure:} Use partial fraction expansion.

For distinct poles $\lambda_1, \lambda_2, \ldots, \lambda_n$:
\begin{equation}
G(s) = \frac{k_1}{s - \lambda_1} + \frac{k_2}{s - \lambda_2} + \cdots + \frac{k_n}{s - \lambda_n}
\end{equation}

where $k_i$ are the residues. The diagonal form is:
\begin{equation}
\mathbf{A}_d = \begin{bmatrix}
\lambda_1 & 0 & \cdots & 0 \\
0 & \lambda_2 & \cdots & 0 \\
\vdots & \vdots & \ddots & \vdots \\
0 & 0 & \cdots & \lambda_n
\end{bmatrix}, \quad
\mathbf{B}_d = \begin{bmatrix} 1 \\ 1 \\ \vdots \\ 1 \end{bmatrix}
\label{eq:diagonal_AB}
\end{equation}

\begin{equation}
\mathbf{C}_d = \begin{bmatrix} k_1 & k_2 & \cdots & k_n \end{bmatrix}, \quad \mathbf{D}_d = 0
\label{eq:diagonal_CD}
\end{equation}

\textbf{Key properties of diagonal form:}
\begin{enumerate}
    \item States are \textbf{decoupled}---each evolves independently according to its eigenvalue.
    \item The \textbf{eigenvalues (poles)} appear directly on the diagonal of $\mathbf{A}_d$.
    \item Stability is immediately visible: stable if all $\text{Re}(\lambda_i) < 0$.
    \item The matrix exponential $e^{\mathbf{A}_d t}$ is trivial to compute.
    \item Only exists for \textbf{distinct eigenvalues} (no repeated poles).
\end{enumerate}

\begin{example}[Diagonal Form via Partial Fractions]
\label{ex:diagonal_form}
Convert to diagonal form:
\begin{equation}
G(s) = \frac{2s + 3}{s^2 + 5s + 6} = \frac{2s + 3}{(s+2)(s+3)}
\end{equation}

\textbf{Step 1: Partial fraction expansion}

\begin{equation}
G(s) = \frac{A}{s+2} + \frac{B}{s+3}
\end{equation}

Using the cover-up method:
\begin{align}
A &= \left.(2s + 3)\right|_{s=-2} \cdot \frac{1}{(-2+3)} = \frac{-4+3}{1} = -1 \\
B &= \left.(2s + 3)\right|_{s=-3} \cdot \frac{1}{(-3+2)} = \frac{-6+3}{-1} = 3
\end{align}

So:
\begin{equation}
G(s) = \frac{-1}{s+2} + \frac{3}{s+3}
\end{equation}

\textbf{Step 2: Construct the diagonal form}

Poles: $\lambda_1 = -2$, $\lambda_2 = -3$. Residues: $k_1 = -1$, $k_2 = 3$.
\begin{equation}
\mathbf{A}_d = \begin{bmatrix} -2 & 0 \\ 0 & -3 \end{bmatrix}, \quad
\mathbf{B}_d = \begin{bmatrix} 1 \\ 1 \end{bmatrix}, \quad
\mathbf{C}_d = \begin{bmatrix} -1 & 3 \end{bmatrix}, \quad
\mathbf{D}_d = 0
\end{equation}

\textbf{Step 3: Verify}

\begin{align}
G(s) &= \mathbf{C}_d(s\mathbf{I} - \mathbf{A}_d)^{-1}\mathbf{B}_d \\
&= \begin{bmatrix} -1 & 3 \end{bmatrix} \begin{bmatrix} \frac{1}{s+2} & 0 \\ 0 & \frac{1}{s+3} \end{bmatrix} \begin{bmatrix} 1 \\ 1 \end{bmatrix} \\
&= \begin{bmatrix} -1 & 3 \end{bmatrix} \begin{bmatrix} \frac{1}{s+2} \\ \frac{1}{s+3} \end{bmatrix} = \frac{-1}{s+2} + \frac{3}{s+3} \quad \checkmark
\end{align}
\end{example}

\subsection{Handling Proper (Non-Strictly-Proper) Transfer Functions}

When the numerator and denominator have the same degree (proper but not strictly proper), there is direct feedthrough:
\begin{equation}
G(s) = \frac{b_n s^n + b_{n-1}s^{n-1} + \cdots + b_0}{s^n + a_{n-1}s^{n-1} + \cdots + a_0}
\end{equation}

\textbf{Procedure:} Perform polynomial long division to extract the constant term.

\begin{example}[Proper Transfer Function with Feedthrough]
Convert to state-space:
\begin{equation}
G(s) = \frac{s^2 + 4s + 3}{s^2 + 2s + 1}
\end{equation}

\textbf{Step 1: Polynomial division}
\begin{equation}
\frac{s^2 + 4s + 3}{s^2 + 2s + 1} = 1 + \frac{2s + 2}{s^2 + 2s + 1}
\end{equation}

\textbf{Step 2: Convert strictly proper part to CCF}

For $G_{sp}(s) = \frac{2s + 2}{s^2 + 2s + 1}$:
\begin{equation}
\mathbf{A} = \begin{bmatrix} 0 & 1 \\ -1 & -2 \end{bmatrix}, \quad
\mathbf{B} = \begin{bmatrix} 0 \\ 1 \end{bmatrix}, \quad
\mathbf{C} = \begin{bmatrix} 2 & 2 \end{bmatrix}
\end{equation}

\textbf{Step 3: Add feedthrough}
\begin{equation}
\mathbf{D} = 1
\end{equation}

The complete state-space representation is $(\mathbf{A}, \mathbf{B}, \mathbf{C}, \mathbf{D})$ with $\mathbf{D} = 1$.
\end{example}

\subsection{Summary: Choosing a Canonical Form}

\begin{center}
\begin{tabular}{llll}
\toprule
\textbf{Form} & \textbf{Guaranteed Property} & \textbf{Best For} & \textbf{Structure} \\
\midrule
Controllable & Controllable & State feedback design & Last row has $-a_i$ \\
Observable & Observable & Observer design & Last column has $-a_i$ \\
Diagonal & Decoupled modes & Modal analysis, stability & Eigenvalues on diagonal \\
Jordan & Handles repeated poles & Non-diagonalizable systems & Block diagonal \\
\bottomrule
\end{tabular}
\end{center}

%=============================================================================
\section{State Transition Matrix and Solution of State Equations}
\label{sec:state_transition_matrix}
%=============================================================================

Having established state-space representations, we now turn to solving the state equations. The solution reveals how the system evolves from initial conditions under applied inputs.

\subsection{Homogeneous Solution: The State Transition Matrix}

Consider the unforced (homogeneous) state equation:
\begin{equation}
\dot{\mathbf{x}}(t) = \mathbf{A}\mathbf{x}(t), \quad \mathbf{x}(0) = \mathbf{x}_0
\end{equation}

For a scalar equation $\dot{x} = ax$, the solution is $x(t) = e^{at}x_0$. The matrix generalization is:
\begin{equation}
\boxed{\mathbf{x}(t) = e^{\mathbf{A}t}\mathbf{x}_0}
\end{equation}

where the \textbf{matrix exponential} $e^{\mathbf{A}t}$ is defined by the power series:
\begin{equation}
e^{\mathbf{A}t} = \mathbf{I} + \mathbf{A}t + \frac{(\mathbf{A}t)^2}{2!} + \frac{(\mathbf{A}t)^3}{3!} + \cdots = \sum_{k=0}^{\infty} \frac{(\mathbf{A}t)^k}{k!}
\label{eq:matrix_exp_series}
\end{equation}

\begin{definition}[State Transition Matrix]
The \textbf{state transition matrix} $\boldsymbol{\Phi}(t) = e^{\mathbf{A}t}$ describes how the state evolves from any initial condition when there is no input.
\end{definition}

\textbf{Properties of the state transition matrix:}
\begin{enumerate}
    \item $\boldsymbol{\Phi}(0) = \mathbf{I}$ (identity at $t=0$)
    \item $\boldsymbol{\Phi}^{-1}(t) = \boldsymbol{\Phi}(-t) = e^{-\mathbf{A}t}$
    \item $\boldsymbol{\Phi}(t_1 + t_2) = \boldsymbol{\Phi}(t_1)\boldsymbol{\Phi}(t_2)$
    \item $\frac{d}{dt}\boldsymbol{\Phi}(t) = \mathbf{A}\boldsymbol{\Phi}(t) = \boldsymbol{\Phi}(t)\mathbf{A}$
    \item $\boldsymbol{\Phi}(t) = \mathcal{L}^{-1}\{(s\mathbf{I} - \mathbf{A})^{-1}\}$
\end{enumerate}

\subsection{Complete Solution with Input}

For the full state equation with input:
\begin{equation}
\dot{\mathbf{x}}(t) = \mathbf{A}\mathbf{x}(t) + \mathbf{B}\mathbf{u}(t), \quad \mathbf{x}(0) = \mathbf{x}_0
\end{equation}

The complete solution is:
\begin{equation}
\boxed{\mathbf{x}(t) = e^{\mathbf{A}t}\mathbf{x}_0 + \int_0^t e^{\mathbf{A}(t-\tau)}\mathbf{B}\mathbf{u}(\tau)\,d\tau}
\label{eq:state_solution}
\end{equation}

This solution has two parts:
\begin{itemize}
    \item \textbf{Free response}: $e^{\mathbf{A}t}\mathbf{x}_0$ (due to initial conditions)
    \item \textbf{Forced response}: $\int_0^t e^{\mathbf{A}(t-\tau)}\mathbf{B}\mathbf{u}(\tau)\,d\tau$ (due to input)
\end{itemize}

The integral is a \textbf{convolution} of the state transition matrix with the input.

\subsection{Computing the Matrix Exponential}

Several methods exist for computing $e^{\mathbf{A}t}$. I will present the most useful approaches.

\subsubsection{Method 1: Laplace Transform}

The state transition matrix can be computed as:
\begin{equation}
e^{\mathbf{A}t} = \mathcal{L}^{-1}\{(s\mathbf{I} - \mathbf{A})^{-1}\}
\label{eq:exp_laplace}
\end{equation}

This method requires computing $(s\mathbf{I} - \mathbf{A})^{-1}$ and taking inverse Laplace transforms element by element.

\begin{example}[Matrix Exponential via Laplace Transform]
\label{ex:matrix_exp_laplace}
Compute $e^{\mathbf{A}t}$ for:
\begin{equation}
\mathbf{A} = \begin{bmatrix} 0 & 1 \\ -2 & -3 \end{bmatrix}
\end{equation}

\textbf{Step 1: Compute $s\mathbf{I} - \mathbf{A}$}
\begin{equation}
s\mathbf{I} - \mathbf{A} = \begin{bmatrix} s & -1 \\ 2 & s+3 \end{bmatrix}
\end{equation}

\textbf{Step 2: Compute the determinant}
\begin{equation}
\det(s\mathbf{I} - \mathbf{A}) = s(s+3) + 2 = s^2 + 3s + 2 = (s+1)(s+2)
\end{equation}

\textbf{Step 3: Compute the inverse}
\begin{equation}
(s\mathbf{I} - \mathbf{A})^{-1} = \frac{1}{(s+1)(s+2)}\begin{bmatrix} s+3 & 1 \\ -2 & s \end{bmatrix}
\end{equation}

\textbf{Step 4: Partial fraction expansion for each element}

For element (1,1): $\frac{s+3}{(s+1)(s+2)} = \frac{2}{s+1} - \frac{1}{s+2}$

For element (1,2): $\frac{1}{(s+1)(s+2)} = \frac{1}{s+1} - \frac{1}{s+2}$

For element (2,1): $\frac{-2}{(s+1)(s+2)} = \frac{-2}{s+1} + \frac{2}{s+2}$

For element (2,2): $\frac{s}{(s+1)(s+2)} = \frac{-1}{s+1} + \frac{2}{s+2}$

\textbf{Step 5: Inverse Laplace transform}
\begin{equation}
e^{\mathbf{A}t} = \begin{bmatrix} 2e^{-t} - e^{-2t} & e^{-t} - e^{-2t} \\ -2e^{-t} + 2e^{-2t} & -e^{-t} + 2e^{-2t} \end{bmatrix}
\end{equation}
\end{example}

\subsubsection{Method 2: Eigenvalue Decomposition (Diagonalizable Case)}

If $\mathbf{A}$ has $n$ linearly independent eigenvectors, we can write:
\begin{equation}
\mathbf{A} = \mathbf{V}\boldsymbol{\Lambda}\mathbf{V}^{-1}
\end{equation}
where $\boldsymbol{\Lambda} = \text{diag}(\lambda_1, \ldots, \lambda_n)$ contains the eigenvalues and $\mathbf{V}$ contains the eigenvectors as columns.

Then:
\begin{equation}
e^{\mathbf{A}t} = \mathbf{V}e^{\boldsymbol{\Lambda}t}\mathbf{V}^{-1} = \mathbf{V}\begin{bmatrix} e^{\lambda_1 t} & & \\ & \ddots & \\ & & e^{\lambda_n t} \end{bmatrix}\mathbf{V}^{-1}
\label{eq:exp_eigen}
\end{equation}

\begin{example}[Matrix Exponential via Eigendecomposition]
\label{ex:matrix_exp_eigen}
For the same matrix from Example~\ref{ex:matrix_exp_laplace}:
\begin{equation}
\mathbf{A} = \begin{bmatrix} 0 & 1 \\ -2 & -3 \end{bmatrix}
\end{equation}

\textbf{Step 1: Find eigenvalues}

From Example~\ref{ex:matrix_exp_laplace}, eigenvalues are $\lambda_1 = -1$, $\lambda_2 = -2$.

\textbf{Step 2: Find eigenvectors}

For $\lambda_1 = -1$: $(\mathbf{A} - (-1)\mathbf{I})\mathbf{v}_1 = \mathbf{0}$
\begin{equation}
\begin{bmatrix} 1 & 1 \\ -2 & -2 \end{bmatrix}\mathbf{v}_1 = \mathbf{0} \implies \mathbf{v}_1 = \begin{bmatrix} 1 \\ -1 \end{bmatrix}
\end{equation}

For $\lambda_2 = -2$: $(\mathbf{A} - (-2)\mathbf{I})\mathbf{v}_2 = \mathbf{0}$
\begin{equation}
\begin{bmatrix} 2 & 1 \\ -2 & -1 \end{bmatrix}\mathbf{v}_2 = \mathbf{0} \implies \mathbf{v}_2 = \begin{bmatrix} 1 \\ -2 \end{bmatrix}
\end{equation}

\textbf{Step 3: Form $\mathbf{V}$ and $\mathbf{V}^{-1}$}
\begin{equation}
\mathbf{V} = \begin{bmatrix} 1 & 1 \\ -1 & -2 \end{bmatrix}, \quad
\mathbf{V}^{-1} = \begin{bmatrix} -2 & -1 \\ 1 & 1 \end{bmatrix}
\end{equation}

\textbf{Step 4: Compute $e^{\mathbf{A}t}$}
\begin{align}
e^{\mathbf{A}t} &= \mathbf{V}\begin{bmatrix} e^{-t} & 0 \\ 0 & e^{-2t} \end{bmatrix}\mathbf{V}^{-1} \\
&= \begin{bmatrix} 1 & 1 \\ -1 & -2 \end{bmatrix}\begin{bmatrix} e^{-t} & 0 \\ 0 & e^{-2t} \end{bmatrix}\begin{bmatrix} -2 & -1 \\ 1 & 1 \end{bmatrix} \\
&= \begin{bmatrix} e^{-t} & e^{-2t} \\ -e^{-t} & -2e^{-2t} \end{bmatrix}\begin{bmatrix} -2 & -1 \\ 1 & 1 \end{bmatrix} \\
&= \begin{bmatrix} -2e^{-t} + e^{-2t} & -e^{-t} + e^{-2t} \\ 2e^{-t} - 2e^{-2t} & e^{-t} - 2e^{-2t} \end{bmatrix}
\end{align}

Wait, this differs in sign from Example~\ref{ex:matrix_exp_laplace}. Let me recheck. Actually, I made an error in the eigenvector calculation. Let me redo it carefully.

For $\lambda_1 = -1$:
\begin{equation}
(\mathbf{A} + \mathbf{I})\mathbf{v}_1 = \begin{bmatrix} 1 & 1 \\ -2 & -2 \end{bmatrix}\mathbf{v}_1 = \mathbf{0}
\end{equation}
This gives $v_{11} + v_{12} = 0$, so $\mathbf{v}_1 = \begin{bmatrix} 1 \\ -1 \end{bmatrix}$.

For $\lambda_2 = -2$:
\begin{equation}
(\mathbf{A} + 2\mathbf{I})\mathbf{v}_2 = \begin{bmatrix} 2 & 1 \\ -2 & -1 \end{bmatrix}\mathbf{v}_2 = \mathbf{0}
\end{equation}
This gives $2v_{21} + v_{22} = 0$, so $\mathbf{v}_2 = \begin{bmatrix} 1 \\ -2 \end{bmatrix}$.

So:
\begin{equation}
\mathbf{V} = \begin{bmatrix} 1 & 1 \\ -1 & -2 \end{bmatrix}
\end{equation}

The determinant is $\det(\mathbf{V}) = -2 - (-1) = -1$, so:
\begin{equation}
\mathbf{V}^{-1} = \frac{1}{-1}\begin{bmatrix} -2 & -1 \\ 1 & 1 \end{bmatrix} = \begin{bmatrix} 2 & 1 \\ -1 & -1 \end{bmatrix}
\end{equation}

Now:
\begin{align}
e^{\mathbf{A}t} &= \begin{bmatrix} 1 & 1 \\ -1 & -2 \end{bmatrix}\begin{bmatrix} e^{-t} & 0 \\ 0 & e^{-2t} \end{bmatrix}\begin{bmatrix} 2 & 1 \\ -1 & -1 \end{bmatrix} \\
&= \begin{bmatrix} e^{-t} & e^{-2t} \\ -e^{-t} & -2e^{-2t} \end{bmatrix}\begin{bmatrix} 2 & 1 \\ -1 & -1 \end{bmatrix} \\
&= \begin{bmatrix} 2e^{-t} - e^{-2t} & e^{-t} - e^{-2t} \\ -2e^{-t} + 2e^{-2t} & -e^{-t} + 2e^{-2t} \end{bmatrix}
\end{align}

This matches Example~\ref{ex:matrix_exp_laplace}.
\end{example}

\subsubsection{Method 3: Diagonal Form (Simplest)}

If the system is already in diagonal form with $\mathbf{A} = \text{diag}(\lambda_1, \ldots, \lambda_n)$:
\begin{equation}
e^{\mathbf{A}t} = \text{diag}(e^{\lambda_1 t}, e^{\lambda_2 t}, \ldots, e^{\lambda_n t})
\end{equation}

This is why diagonal canonical form is valuable for analysis---the matrix exponential is trivial to compute.

\subsection{Step Response via State-Space}

Let us compute the step response using the state-space solution.

\begin{example}[Step Response of Second-Order System]
\label{ex:step_response_ss}
For the controllable canonical form from Example~\ref{ex:tf_to_ccf_2nd}:
\begin{equation}
\mathbf{A} = \begin{bmatrix} 0 & 1 \\ -6 & -5 \end{bmatrix}, \quad
\mathbf{B} = \begin{bmatrix} 0 \\ 1 \end{bmatrix}, \quad
\mathbf{C} = \begin{bmatrix} 3 & 2 \end{bmatrix}, \quad
\mathbf{D} = 0
\end{equation}

With unit step input $u(t) = 1$ and zero initial conditions $\mathbf{x}_0 = \mathbf{0}$.

\textbf{Step 1: Compute the forced response}

From equation~\eqref{eq:state_solution}:
\begin{equation}
\mathbf{x}(t) = \int_0^t e^{\mathbf{A}(t-\tau)}\mathbf{B}\,d\tau = \int_0^t e^{\mathbf{A}\sigma}\mathbf{B}\,d\sigma
\end{equation}
where I substituted $\sigma = t - \tau$.

We found that the eigenvalues are $\lambda_1 = -2$, $\lambda_2 = -3$. Using the diagonal form from Example~\ref{ex:diagonal_form}:
\begin{equation}
\mathbf{A}_d = \begin{bmatrix} -2 & 0 \\ 0 & -3 \end{bmatrix}, \quad e^{\mathbf{A}_d t} = \begin{bmatrix} e^{-2t} & 0 \\ 0 & e^{-3t} \end{bmatrix}
\end{equation}

The state in diagonal coordinates:
\begin{equation}
\mathbf{z}(t) = \int_0^t \begin{bmatrix} e^{-2\sigma} & 0 \\ 0 & e^{-3\sigma} \end{bmatrix}\begin{bmatrix} 1 \\ 1 \end{bmatrix}d\sigma = \begin{bmatrix} \frac{1 - e^{-2t}}{2} \\ \frac{1 - e^{-3t}}{3} \end{bmatrix}
\end{equation}

\textbf{Step 2: Compute the output}
\begin{align}
y(t) &= \mathbf{C}_d\mathbf{z}(t) = \begin{bmatrix} -1 & 3 \end{bmatrix}\begin{bmatrix} \frac{1 - e^{-2t}}{2} \\ \frac{1 - e^{-3t}}{3} \end{bmatrix} \\
&= -\frac{1 - e^{-2t}}{2} + \frac{3(1 - e^{-3t})}{3} \\
&= -\frac{1}{2} + \frac{e^{-2t}}{2} + 1 - e^{-3t} \\
&= \frac{1}{2} + \frac{1}{2}e^{-2t} - e^{-3t}
\end{align}

\textbf{Verification:} The steady-state value is $y(\infty) = \frac{1}{2}$, which equals $G(0) = \frac{3}{6} = \frac{1}{2}$. $\checkmark$
\end{example}

%=============================================================================
\section{Controllability}
\label{sec:controllability}
%=============================================================================

Controllability is one of the most fundamental concepts in modern control theory. It answers a critical question: \textit{Can we steer the system from any initial state to any desired final state using the available inputs?}

\subsection{Definition of Controllability}

\begin{definition}[Controllability]
A linear time-invariant system $(\mathbf{A}, \mathbf{B})$ is \textbf{controllable} if, for any initial state $\mathbf{x}_0$ and any desired final state $\mathbf{x}_f$, there exists an input $\mathbf{u}(t)$ over a finite time interval $[0, t_f]$ that drives the system from $\mathbf{x}_0$ to $\mathbf{x}_f$.
\end{definition}

Controllability is a property of the pair $(\mathbf{A}, \mathbf{B})$---it depends on the system dynamics and how inputs enter the system. The output matrices $\mathbf{C}$ and $\mathbf{D}$ do not affect controllability.

\subsection{The Controllability Matrix}

The key to testing controllability is the \textbf{controllability matrix}:
\begin{equation}
\boxed{\mathcal{C} = \begin{bmatrix} \mathbf{B} & \mathbf{A}\mathbf{B} & \mathbf{A}^2\mathbf{B} & \cdots & \mathbf{A}^{n-1}\mathbf{B} \end{bmatrix}}
\label{eq:controllability_matrix}
\end{equation}

For a system with $n$ states and $m$ inputs, $\mathcal{C}$ is an $n \times nm$ matrix.

\begin{theorem}[Controllability Rank Test]
\label{thm:controllability_rank}
The system $(\mathbf{A}, \mathbf{B})$ is controllable if and only if:
\begin{equation}
\text{rank}(\mathcal{C}) = n
\end{equation}
\end{theorem}

\textbf{Physical interpretation:}
\begin{itemize}
    \item $\mathbf{B}$: States directly reachable from input in one time step
    \item $\mathbf{A}\mathbf{B}$: States reachable after the dynamics propagate the input effect once
    \item $\mathbf{A}^2\mathbf{B}$: States reachable after two propagation steps
    \item $\mathbf{A}^{n-1}\mathbf{B}$: States reachable after $n-1$ propagation steps
\end{itemize}

If rank$(\mathcal{C}) = n$, every state direction can be reached through some combination of these influences.

\begin{example}[Controllability Test for Mass-Spring-Damper]
\label{ex:controllability_msd}
Test controllability of the mass-spring-damper system with $m = 1$, $k = 4$, $b = 2$:
\begin{equation}
\mathbf{A} = \begin{bmatrix} 0 & 1 \\ -4 & -2 \end{bmatrix}, \quad \mathbf{B} = \begin{bmatrix} 0 \\ 1 \end{bmatrix}
\end{equation}

\textbf{Step 1: Compute $\mathbf{A}\mathbf{B}$}
\begin{equation}
\mathbf{A}\mathbf{B} = \begin{bmatrix} 0 & 1 \\ -4 & -2 \end{bmatrix}\begin{bmatrix} 0 \\ 1 \end{bmatrix} = \begin{bmatrix} 1 \\ -2 \end{bmatrix}
\end{equation}

\textbf{Step 2: Form the controllability matrix}
\begin{equation}
\mathcal{C} = \begin{bmatrix} \mathbf{B} & \mathbf{A}\mathbf{B} \end{bmatrix} = \begin{bmatrix} 0 & 1 \\ 1 & -2 \end{bmatrix}
\end{equation}

\textbf{Step 3: Check rank}
\begin{equation}
\det(\mathcal{C}) = 0 \cdot (-2) - 1 \cdot 1 = -1 \neq 0
\end{equation}

Since $\det(\mathcal{C}) \neq 0$, rank$(\mathcal{C}) = 2 = n$. The system is \textbf{controllable}.

\textbf{Physical interpretation:} Even though the force input does not directly affect position (only velocity), the coupling through the dynamics allows us to control both position and velocity.
\end{example}

\begin{example}[Uncontrollable System]
\label{ex:uncontrollable}
Consider:
\begin{equation}
\mathbf{A} = \begin{bmatrix} -1 & 0 \\ 0 & -2 \end{bmatrix}, \quad \mathbf{B} = \begin{bmatrix} 1 \\ 0 \end{bmatrix}
\end{equation}

\textbf{Step 1: Compute $\mathbf{A}\mathbf{B}$}
\begin{equation}
\mathbf{A}\mathbf{B} = \begin{bmatrix} -1 & 0 \\ 0 & -2 \end{bmatrix}\begin{bmatrix} 1 \\ 0 \end{bmatrix} = \begin{bmatrix} -1 \\ 0 \end{bmatrix}
\end{equation}

\textbf{Step 2: Form the controllability matrix}
\begin{equation}
\mathcal{C} = \begin{bmatrix} 1 & -1 \\ 0 & 0 \end{bmatrix}
\end{equation}

\textbf{Step 3: Check rank}

The second row is all zeros, so rank$(\mathcal{C}) = 1 < 2$. The system is \textbf{not controllable}.

\textbf{Physical interpretation:} The input only affects $x_1$. Since $\mathbf{A}$ is diagonal, there is no coupling, so the input can never influence $x_2$. State $x_2$ evolves according to $\dot{x}_2 = -2x_2$ regardless of the input.
\end{example}

\subsection{The Controllability Gramian}

An alternative approach to analyzing controllability uses the \textbf{controllability Gramian}.

\begin{definition}[Controllability Gramian]
For a stable system ($\mathbf{A}$ has all eigenvalues with negative real parts), the controllability Gramian is:
\begin{equation}
\mathbf{W}_c = \int_0^\infty e^{\mathbf{A}t}\mathbf{B}\mathbf{B}^\top e^{\mathbf{A}^\top t}\,dt
\label{eq:controllability_gramian}
\end{equation}
\end{definition}

The Gramian satisfies the \textbf{Lyapunov equation}:
\begin{equation}
\mathbf{A}\mathbf{W}_c + \mathbf{W}_c\mathbf{A}^\top + \mathbf{B}\mathbf{B}^\top = \mathbf{0}
\label{eq:lyapunov_controllability}
\end{equation}

\begin{theorem}[Controllability via Gramian]
A stable system $(\mathbf{A}, \mathbf{B})$ is controllable if and only if $\mathbf{W}_c$ is positive definite (i.e., $\mathbf{W}_c > 0$).
\end{theorem}

The Gramian approach is computationally advantageous for large systems because solving the Lyapunov equation~\eqref{eq:lyapunov_controllability} is efficient ($O(n^3)$ operations) compared to computing and checking the rank of $\mathcal{C}$.

\textbf{Physical interpretation of the Gramian:} The eigenvalues of $\mathbf{W}_c$ indicate how \textit{easy} it is to control different state directions. Large eigenvalues correspond to directions that are easy to control; small eigenvalues correspond to directions that require large control effort.

\subsection{Controllability of Canonical Forms}

The controllability of canonical forms follows directly from their structure:

\begin{theorem}[Controllability of CCF]
The controllable canonical form $(\mathbf{A}_c, \mathbf{B}_c)$ is always controllable. The controllability matrix has determinant $\pm 1$.
\end{theorem}

\begin{proof}
For a second-order CCF:
\begin{equation}
\mathbf{A}_c = \begin{bmatrix} 0 & 1 \\ -a_0 & -a_1 \end{bmatrix}, \quad \mathbf{B}_c = \begin{bmatrix} 0 \\ 1 \end{bmatrix}
\end{equation}

The controllability matrix is:
\begin{equation}
\mathcal{C} = \begin{bmatrix} 0 & 1 \\ 1 & -a_1 \end{bmatrix}
\end{equation}

We have $\det(\mathcal{C}) = -1 \neq 0$, so rank$(\mathcal{C}) = 2$. This generalizes to arbitrary order.
\end{proof}

\textbf{Warning:} Observable canonical form is \textit{not} guaranteed to be controllable, and diagonal form may or may not be controllable depending on the $\mathbf{B}$ matrix structure.

%=============================================================================
\section{Observability}
\label{sec:observability}
%=============================================================================

Observability is the dual concept to controllability. It answers: \textit{Can we determine the internal state of the system from measurements of the output?}

\subsection{Definition of Observability}

\begin{definition}[Observability]
A linear time-invariant system $(\mathbf{A}, \mathbf{C})$ is \textbf{observable} if the initial state $\mathbf{x}_0$ can be uniquely determined from knowledge of the input $\mathbf{u}(t)$ and output $\mathbf{y}(t)$ over a finite time interval $[0, t_f]$.
\end{definition}

Observability is a property of the pair $(\mathbf{A}, \mathbf{C})$---it depends on the system dynamics and what is measured. The input matrices $\mathbf{B}$ and $\mathbf{D}$ do not affect observability.

\subsection{The Observability Matrix}

\begin{equation}
\boxed{\mathcal{O} = \begin{bmatrix} \mathbf{C} \\ \mathbf{C}\mathbf{A} \\ \mathbf{C}\mathbf{A}^2 \\ \vdots \\ \mathbf{C}\mathbf{A}^{n-1} \end{bmatrix}}
\label{eq:observability_matrix}
\end{equation}

For a system with $n$ states and $p$ outputs, $\mathcal{O}$ is a $pn \times n$ matrix.

\begin{theorem}[Observability Rank Test]
\label{thm:observability_rank}
The system $(\mathbf{A}, \mathbf{C})$ is observable if and only if:
\begin{equation}
\text{rank}(\mathcal{O}) = n
\end{equation}
\end{theorem}

\textbf{Physical interpretation:}
\begin{itemize}
    \item $\mathbf{C}$: States directly visible in the output
    \item $\mathbf{C}\mathbf{A}$: States whose effects appear in the output after one time step
    \item $\mathbf{C}\mathbf{A}^{n-1}$: States whose effects appear after $n-1$ time steps
\end{itemize}

If rank$(\mathcal{O}) = n$, every state can eventually be distinguished through its effect on the output.

\begin{example}[Observability Test]
\label{ex:observability_test}
Test observability of:
\begin{equation}
\mathbf{A} = \begin{bmatrix} 0 & 1 \\ -6 & -5 \end{bmatrix}, \quad \mathbf{C} = \begin{bmatrix} 1 & 0 \end{bmatrix}
\end{equation}

This represents a system where only position is measured (not velocity).

\textbf{Step 1: Compute $\mathbf{C}\mathbf{A}$}
\begin{equation}
\mathbf{C}\mathbf{A} = \begin{bmatrix} 1 & 0 \end{bmatrix}\begin{bmatrix} 0 & 1 \\ -6 & -5 \end{bmatrix} = \begin{bmatrix} 0 & 1 \end{bmatrix}
\end{equation}

\textbf{Step 2: Form the observability matrix}
\begin{equation}
\mathcal{O} = \begin{bmatrix} \mathbf{C} \\ \mathbf{C}\mathbf{A} \end{bmatrix} = \begin{bmatrix} 1 & 0 \\ 0 & 1 \end{bmatrix}
\end{equation}

\textbf{Step 3: Check rank}

$\mathcal{O}$ is the identity matrix, so rank$(\mathcal{O}) = 2 = n$. The system is \textbf{observable}.

\textbf{Physical interpretation:} Although we only measure position directly, velocity affects how position changes. By observing position over time, we can infer velocity.
\end{example}

\begin{example}[Unobservable System]
\label{ex:unobservable}
Consider:
\begin{equation}
\mathbf{A} = \begin{bmatrix} -1 & 0 \\ 0 & -2 \end{bmatrix}, \quad \mathbf{C} = \begin{bmatrix} 1 & 0 \end{bmatrix}
\end{equation}

\textbf{Step 1: Compute $\mathbf{C}\mathbf{A}$}
\begin{equation}
\mathbf{C}\mathbf{A} = \begin{bmatrix} 1 & 0 \end{bmatrix}\begin{bmatrix} -1 & 0 \\ 0 & -2 \end{bmatrix} = \begin{bmatrix} -1 & 0 \end{bmatrix}
\end{equation}

\textbf{Step 2: Form the observability matrix}
\begin{equation}
\mathcal{O} = \begin{bmatrix} 1 & 0 \\ -1 & 0 \end{bmatrix}
\end{equation}

\textbf{Step 3: Check rank}

The second column is all zeros, so rank$(\mathcal{O}) = 1 < 2$. The system is \textbf{not observable}.

\textbf{Physical interpretation:} We only measure $x_1$. Since $\mathbf{A}$ is diagonal (no coupling), $x_2$ has no effect on $x_1$ or the output. We cannot determine $x_2$ from measurements of $y = x_1$.
\end{example}

\subsection{The Observability Gramian}

\begin{definition}[Observability Gramian]
For a stable system, the observability Gramian is:
\begin{equation}
\mathbf{W}_o = \int_0^\infty e^{\mathbf{A}^\top t}\mathbf{C}^\top\mathbf{C} e^{\mathbf{A}t}\,dt
\label{eq:observability_gramian}
\end{equation}
\end{definition}

The Gramian satisfies:
\begin{equation}
\mathbf{A}^\top\mathbf{W}_o + \mathbf{W}_o\mathbf{A} + \mathbf{C}^\top\mathbf{C} = \mathbf{0}
\label{eq:lyapunov_observability}
\end{equation}

\begin{theorem}[Observability via Gramian]
A stable system $(\mathbf{A}, \mathbf{C})$ is observable if and only if $\mathbf{W}_o > 0$ (positive definite).
\end{theorem}

\textbf{Physical interpretation:} The eigenvalues of $\mathbf{W}_o$ indicate how \textit{easy} it is to observe different state directions. Large eigenvalues correspond to states that strongly affect the output; small eigenvalues correspond to states that are difficult to observe.

\subsection{Duality Between Controllability and Observability}

There is a beautiful duality between controllability and observability:

\begin{theorem}[Duality]
\label{thm:duality}
The system $(\mathbf{A}, \mathbf{B})$ is controllable if and only if the dual system $(\mathbf{A}^\top, \mathbf{B}^\top)$ is observable.

Equivalently, $(\mathbf{A}, \mathbf{C})$ is observable if and only if $(\mathbf{A}^\top, \mathbf{C}^\top)$ is controllable.
\end{theorem}

This duality means:
\begin{itemize}
    \item Every theorem about controllability has a corresponding theorem about observability
    \item The controllability matrix of $(\mathbf{A}, \mathbf{B})$ equals the transpose of the observability matrix of $(\mathbf{A}^\top, \mathbf{B}^\top)$
    \item Observer design for $(\mathbf{A}, \mathbf{C})$ is mathematically equivalent to controller design for $(\mathbf{A}^\top, \mathbf{C}^\top)$
\end{itemize}

\textbf{This is why observable canonical form is the transpose of controllable canonical form.}

\subsection{Minimal Realization}

\begin{definition}[Minimal Realization]
A state-space realization $(\mathbf{A}, \mathbf{B}, \mathbf{C}, \mathbf{D})$ is \textbf{minimal} if it has the smallest possible number of states among all realizations of the same transfer function.
\end{definition}

\begin{theorem}[Minimal Realization Criterion]
A realization is minimal if and only if it is both controllable and observable.
\end{theorem}

\textbf{Implications:}
\begin{itemize}
    \item A minimal realization has $n$ states, where $n$ equals the degree of the denominator of the transfer function (with no pole-zero cancellations)
    \item Non-minimal realizations have \textbf{hidden modes}---uncontrollable or unobservable states that do not appear in the transfer function
    \item Converting a transfer function to CCF or OCF produces a minimal realization (assuming no pole-zero cancellations in the original transfer function)
\end{itemize}

\begin{example}[Non-Minimal Realization]
\label{ex:non_minimal}
Consider:
\begin{equation}
\mathbf{A} = \begin{bmatrix} -1 & 0 \\ 0 & -2 \end{bmatrix}, \quad \mathbf{B} = \begin{bmatrix} 1 \\ 0 \end{bmatrix}, \quad \mathbf{C} = \begin{bmatrix} 1 & 0 \end{bmatrix}, \quad \mathbf{D} = 0
\end{equation}

From Examples~\ref{ex:uncontrollable} and~\ref{ex:unobservable}, this system is neither controllable nor observable.

The transfer function is:
\begin{equation}
G(s) = \mathbf{C}(s\mathbf{I} - \mathbf{A})^{-1}\mathbf{B} = \begin{bmatrix} 1 & 0 \end{bmatrix}\begin{bmatrix} \frac{1}{s+1} & 0 \\ 0 & \frac{1}{s+2} \end{bmatrix}\begin{bmatrix} 1 \\ 0 \end{bmatrix} = \frac{1}{s+1}
\end{equation}

The pole at $s = -2$ does not appear in the transfer function! This is a \textbf{hidden mode}---the state $x_2$ evolves according to $\dot{x}_2 = -2x_2$ but cannot be controlled or observed.

A minimal realization is simply:
\begin{equation}
A_{\min} = -1, \quad B_{\min} = 1, \quad C_{\min} = 1, \quad D_{\min} = 0
\end{equation}
This first-order system has the same transfer function but no hidden modes.
\end{example}

%=============================================================================
\section{Canonical Forms}
\label{sec:canonical_forms}
%=============================================================================

We have already introduced controllable, observable, and diagonal canonical forms in Section~\ref{sec:tf_to_ss}. Here, I provide a comprehensive treatment including the Jordan canonical form for systems with repeated eigenvalues.

\subsection{Summary of Canonical Forms}

The following table summarizes the key canonical forms:

\begin{center}
\begin{tabular}{p{2.5cm}p{4cm}p{4cm}p{3cm}}
\toprule
\textbf{Form} & \textbf{A Matrix Structure} & \textbf{Key Feature} & \textbf{Best Application} \\
\midrule
Controllable & Companion (last row has $-a_i$) & Always controllable & State feedback design \\
Observable & Companion (last col has $-a_i$) & Always observable & Observer design \\
Diagonal & Diagonal (eigenvalues) & Decoupled modes & Modal analysis \\
Jordan & Block diagonal & Handles repeated poles & General analysis \\
\bottomrule
\end{tabular}
\end{center}

\subsection{Jordan Canonical Form}

When the system matrix $\mathbf{A}$ has \textbf{repeated eigenvalues}, it may not be diagonalizable. In this case, we use the \textbf{Jordan canonical form}.

\begin{definition}[Jordan Block]
A Jordan block $\mathbf{J}_k(\lambda)$ of size $k$ associated with eigenvalue $\lambda$ is:
\begin{equation}
\mathbf{J}_k(\lambda) = \begin{bmatrix}
\lambda & 1 & 0 & \cdots & 0 \\
0 & \lambda & 1 & \cdots & 0 \\
\vdots & & \ddots & \ddots & \vdots \\
0 & 0 & \cdots & \lambda & 1 \\
0 & 0 & \cdots & 0 & \lambda
\end{bmatrix}_{k \times k}
\end{equation}
\end{definition}

The Jordan canonical form of $\mathbf{A}$ is a block diagonal matrix:
\begin{equation}
\mathbf{J} = \begin{bmatrix}
\mathbf{J}_{k_1}(\lambda_1) & & \\
& \mathbf{J}_{k_2}(\lambda_2) & \\
& & \ddots
\end{bmatrix}
\end{equation}

\begin{example}[Jordan Form with Repeated Eigenvalue]
\label{ex:jordan_form}
Consider a system with double pole at $s = -2$:
\begin{equation}
G(s) = \frac{1}{(s+2)^2}
\end{equation}

The Jordan form realization is:
\begin{equation}
\mathbf{A}_J = \begin{bmatrix} -2 & 1 \\ 0 & -2 \end{bmatrix}, \quad
\mathbf{B}_J = \begin{bmatrix} 0 \\ 1 \end{bmatrix}, \quad
\mathbf{C}_J = \begin{bmatrix} 1 & 0 \end{bmatrix}, \quad
\mathbf{D}_J = 0
\end{equation}

The ``1'' above the diagonal in $\mathbf{A}_J$ couples the two states---this is essential because with a repeated eigenvalue, we cannot have fully decoupled dynamics.

\textbf{Matrix exponential for Jordan blocks:}
\begin{equation}
e^{\mathbf{J}_2(\lambda)t} = e^{\lambda t}\begin{bmatrix} 1 & t \\ 0 & 1 \end{bmatrix}
\end{equation}

For our example:
\begin{equation}
e^{\mathbf{A}_J t} = e^{-2t}\begin{bmatrix} 1 & t \\ 0 & 1 \end{bmatrix}
\end{equation}

Note the polynomial term $t$ multiplying the exponential---this is characteristic of repeated poles.
\end{example}

\subsection{Relationship Between Canonical Forms}

All canonical forms represent the \textit{same} input-output behavior (same transfer function). They differ only in:
\begin{itemize}
    \item How states are defined (different coordinate systems)
    \item Which structural property is emphasized (controllability, observability, modal decomposition)
    \item Numerical conditioning for computation
\end{itemize}

Any two minimal realizations of the same transfer function are related by a similarity transformation.

%=============================================================================
\section{Similarity Transformations}
\label{sec:similarity_transformations}
%=============================================================================

Similarity transformations are the mathematical tool for converting between different state-space realizations. They represent a change of coordinates in state space.

\subsection{Definition and Properties}

\begin{definition}[Similarity Transformation]
Given a state-space system $(\mathbf{A}, \mathbf{B}, \mathbf{C}, \mathbf{D})$ and an invertible matrix $\mathbf{T} \in \mathbb{R}^{n \times n}$, the \textbf{similarity transformation} $\bar{\mathbf{x}} = \mathbf{T}\mathbf{x}$ produces a new realization:
\begin{align}
\bar{\mathbf{A}} &= \mathbf{T}\mathbf{A}\mathbf{T}^{-1} \\
\bar{\mathbf{B}} &= \mathbf{T}\mathbf{B} \\
\bar{\mathbf{C}} &= \mathbf{C}\mathbf{T}^{-1} \\
\bar{\mathbf{D}} &= \mathbf{D}
\end{align}
\end{definition}

\begin{theorem}[Invariants Under Similarity Transformation]
The following properties are preserved under similarity transformations:
\begin{enumerate}
    \item \textbf{Transfer function:} $\bar{G}(s) = G(s)$
    \item \textbf{Eigenvalues:} $\sigma(\bar{\mathbf{A}}) = \sigma(\mathbf{A})$ (same poles)
    \item \textbf{Controllability:} rank$(\bar{\mathcal{C}})$ = rank$(\mathcal{C})$
    \item \textbf{Observability:} rank$(\bar{\mathcal{O}})$ = rank$(\mathcal{O})$
    \item \textbf{Characteristic polynomial:} $\det(s\mathbf{I} - \bar{\mathbf{A}}) = \det(s\mathbf{I} - \mathbf{A})$
\end{enumerate}
\end{theorem}

\begin{proof}[Proof of transfer function invariance]
\begin{align}
\bar{G}(s) &= \bar{\mathbf{C}}(s\mathbf{I} - \bar{\mathbf{A}})^{-1}\bar{\mathbf{B}} + \bar{\mathbf{D}} \\
&= \mathbf{C}\mathbf{T}^{-1}(s\mathbf{I} - \mathbf{T}\mathbf{A}\mathbf{T}^{-1})^{-1}\mathbf{T}\mathbf{B} + \mathbf{D} \\
&= \mathbf{C}\mathbf{T}^{-1}[\mathbf{T}(s\mathbf{I} - \mathbf{A})\mathbf{T}^{-1}]^{-1}\mathbf{T}\mathbf{B} + \mathbf{D} \\
&= \mathbf{C}\mathbf{T}^{-1}\mathbf{T}(s\mathbf{I} - \mathbf{A})^{-1}\mathbf{T}^{-1}\mathbf{T}\mathbf{B} + \mathbf{D} \\
&= \mathbf{C}(s\mathbf{I} - \mathbf{A})^{-1}\mathbf{B} + \mathbf{D} = G(s)
\end{align}
\end{proof}

\subsection{Finding the Transformation Matrix}

Given two realizations $(\mathbf{A}_1, \mathbf{B}_1, \mathbf{C}_1)$ and $(\mathbf{A}_2, \mathbf{B}_2, \mathbf{C}_2)$ of the same transfer function, we want to find $\mathbf{T}$ such that:
\begin{equation}
\mathbf{A}_2 = \mathbf{T}\mathbf{A}_1\mathbf{T}^{-1}, \quad \mathbf{B}_2 = \mathbf{T}\mathbf{B}_1, \quad \mathbf{C}_2 = \mathbf{C}_1\mathbf{T}^{-1}
\end{equation}

\subsubsection{Method 1: Using Controllability Matrices}

If both realizations are controllable:
\begin{equation}
\mathbf{T} = \mathcal{C}_2 \mathcal{C}_1^{-1}
\label{eq:T_controllability}
\end{equation}
where $\mathcal{C}_1$ and $\mathcal{C}_2$ are the controllability matrices.

\subsubsection{Method 2: Using Observability Matrices}

If both realizations are observable:
\begin{equation}
\mathbf{T} = \mathcal{O}_1^{-1}\mathcal{O}_2
\label{eq:T_observability}
\end{equation}
where $\mathcal{O}_1$ and $\mathcal{O}_2$ are the observability matrices.

\subsubsection{Method 3: Eigendecomposition (for Diagonal Form)}

To transform to diagonal form:
\begin{equation}
\mathbf{T} = \mathbf{V}^{-1}
\end{equation}
where $\mathbf{V}$ is the eigenvector matrix of $\mathbf{A}$ (columns are eigenvectors).

\begin{example}[Transformation from CCF to Diagonal Form]
\label{ex:ccf_to_diagonal}
Transform the controllable canonical form:
\begin{equation}
\mathbf{A}_c = \begin{bmatrix} 0 & 1 \\ -6 & -5 \end{bmatrix}, \quad \mathbf{B}_c = \begin{bmatrix} 0 \\ 1 \end{bmatrix}, \quad \mathbf{C}_c = \begin{bmatrix} 3 & 2 \end{bmatrix}
\end{equation}
to diagonal form.

\textbf{Step 1: Find eigenvalues of $\mathbf{A}_c$}
\begin{equation}
\det(\lambda\mathbf{I} - \mathbf{A}_c) = \lambda^2 + 5\lambda + 6 = (\lambda + 2)(\lambda + 3) = 0
\end{equation}
Eigenvalues: $\lambda_1 = -2$, $\lambda_2 = -3$.

\textbf{Step 2: Find eigenvectors}

For $\lambda_1 = -2$:
\begin{equation}
(\mathbf{A}_c + 2\mathbf{I})\mathbf{v}_1 = \begin{bmatrix} 2 & 1 \\ -6 & -3 \end{bmatrix}\mathbf{v}_1 = \mathbf{0} \implies \mathbf{v}_1 = \begin{bmatrix} 1 \\ -2 \end{bmatrix}
\end{equation}

For $\lambda_2 = -3$:
\begin{equation}
(\mathbf{A}_c + 3\mathbf{I})\mathbf{v}_2 = \begin{bmatrix} 3 & 1 \\ -6 & -2 \end{bmatrix}\mathbf{v}_2 = \mathbf{0} \implies \mathbf{v}_2 = \begin{bmatrix} 1 \\ -3 \end{bmatrix}
\end{equation}

\textbf{Step 3: Form transformation matrix}
\begin{equation}
\mathbf{V} = \begin{bmatrix} 1 & 1 \\ -2 & -3 \end{bmatrix}, \quad \mathbf{T} = \mathbf{V}^{-1} = \begin{bmatrix} 3 & 1 \\ -2 & -1 \end{bmatrix}
\end{equation}

\textbf{Step 4: Transform the system}
\begin{align}
\mathbf{A}_d &= \mathbf{T}\mathbf{A}_c\mathbf{T}^{-1} = \begin{bmatrix} -2 & 0 \\ 0 & -3 \end{bmatrix} \\
\mathbf{B}_d &= \mathbf{T}\mathbf{B}_c = \begin{bmatrix} 3 & 1 \\ -2 & -1 \end{bmatrix}\begin{bmatrix} 0 \\ 1 \end{bmatrix} = \begin{bmatrix} 1 \\ -1 \end{bmatrix} \\
\mathbf{C}_d &= \mathbf{C}_c\mathbf{T}^{-1} = \begin{bmatrix} 3 & 2 \end{bmatrix}\begin{bmatrix} 1 & 1 \\ -2 & -3 \end{bmatrix} = \begin{bmatrix} -1 & -3 \end{bmatrix}
\end{align}

\textbf{Verification:} The diagonal form has eigenvalues $-2$ and $-3$ on the diagonal, matching the poles. The transfer function is:
\begin{align}
G(s) &= \mathbf{C}_d(s\mathbf{I} - \mathbf{A}_d)^{-1}\mathbf{B}_d \\
&= \begin{bmatrix} -1 & -3 \end{bmatrix}\begin{bmatrix} \frac{1}{s+2} & 0 \\ 0 & \frac{1}{s+3} \end{bmatrix}\begin{bmatrix} 1 \\ -1 \end{bmatrix} \\
&= \frac{-1}{s+2} + \frac{3}{s+3} = \frac{2s+3}{(s+2)(s+3)} \quad \checkmark
\end{align}
\end{example}

\subsection{Transformation to Controllable Canonical Form}

A particularly useful transformation converts any controllable realization to controllable canonical form.

\begin{theorem}[Transformation to CCF]
\label{thm:transform_to_ccf}
Let $(\mathbf{A}, \mathbf{B})$ be controllable with controllability matrix $\mathcal{C}$. The transformation to controllable canonical form is:
\begin{equation}
\mathbf{T} = \mathcal{C}_c \mathcal{C}^{-1}
\end{equation}
where $\mathcal{C}_c$ is the controllability matrix of the CCF.

Alternatively, define:
\begin{equation}
\mathbf{T} = \begin{bmatrix} \mathbf{q}^\top \\ \mathbf{q}^\top\mathbf{A} \\ \vdots \\ \mathbf{q}^\top\mathbf{A}^{n-1} \end{bmatrix}
\end{equation}
where $\mathbf{q}^\top$ is the last row of $\mathcal{C}^{-1}$.
\end{theorem}

%=============================================================================
\section{Practical Applications and Worked Examples}
\label{sec:practical_applications}
%=============================================================================

Now I will present comprehensive worked examples that integrate the concepts from this chapter. These examples progress from fundamental conversions to realistic engineering systems.

\subsection{Example: Deriving State-Space from a Physical System}

\begin{example}[DC Motor State-Space Model]
\label{ex:dc_motor}
Derive the state-space model for a DC motor with:
\begin{itemize}
    \item Armature resistance $R_a$, inductance $L_a$
    \item Motor torque constant $K_t$, back-EMF constant $K_e$
    \item Rotor inertia $J$, viscous friction $b$
    \item Input: armature voltage $v_a(t)$
    \item Outputs: angular velocity $\omega(t)$ and armature current $i_a(t)$
\end{itemize}

\textbf{Step 1: Write the governing equations}

\textit{Electrical subsystem:}
\begin{equation}
L_a \frac{di_a}{dt} = v_a - R_a i_a - K_e \omega
\label{eq:dc_motor_elec}
\end{equation}

\textit{Mechanical subsystem:}
\begin{equation}
J\frac{d\omega}{dt} = K_t i_a - b\omega
\label{eq:dc_motor_mech}
\end{equation}

\textbf{Step 2: Define state variables}

Energy storage elements:
\begin{itemize}
    \item Inductor $L_a$: stores magnetic energy $\frac{1}{2}L_a i_a^2$ $\Rightarrow$ state $x_1 = i_a$
    \item Inertia $J$: stores kinetic energy $\frac{1}{2}J\omega^2$ $\Rightarrow$ state $x_2 = \omega$
\end{itemize}

\textbf{Step 3: Write state equations}

From equation~\eqref{eq:dc_motor_elec}:
\begin{equation}
\dot{x}_1 = \frac{di_a}{dt} = -\frac{R_a}{L_a}i_a - \frac{K_e}{L_a}\omega + \frac{1}{L_a}v_a = -\frac{R_a}{L_a}x_1 - \frac{K_e}{L_a}x_2 + \frac{1}{L_a}v_a
\end{equation}

From equation~\eqref{eq:dc_motor_mech}:
\begin{equation}
\dot{x}_2 = \frac{d\omega}{dt} = \frac{K_t}{J}i_a - \frac{b}{J}\omega = \frac{K_t}{J}x_1 - \frac{b}{J}x_2
\end{equation}

\textbf{Step 4: Assemble the state-space model}
\begin{equation}
\begin{bmatrix} \dot{x}_1 \\ \dot{x}_2 \end{bmatrix} = \begin{bmatrix} -\frac{R_a}{L_a} & -\frac{K_e}{L_a} \\ \frac{K_t}{J} & -\frac{b}{J} \end{bmatrix} \begin{bmatrix} x_1 \\ x_2 \end{bmatrix} + \begin{bmatrix} \frac{1}{L_a} \\ 0 \end{bmatrix} v_a
\end{equation}

With both outputs measured:
\begin{equation}
\begin{bmatrix} y_1 \\ y_2 \end{bmatrix} = \begin{bmatrix} i_a \\ \omega \end{bmatrix} = \begin{bmatrix} 1 & 0 \\ 0 & 1 \end{bmatrix} \begin{bmatrix} x_1 \\ x_2 \end{bmatrix}
\end{equation}

\textbf{Step 5: Numerical example}

Let $R_a = 1\,\Omega$, $L_a = 0.5$\,H, $K_t = K_e = 0.1$\,V$\cdot$s/rad, $J = 0.01$\,kg$\cdot$m$^2$, $b = 0.1$\,N$\cdot$m$\cdot$s/rad:
\begin{equation}
\mathbf{A} = \begin{bmatrix} -2 & -0.2 \\ 10 & -10 \end{bmatrix}, \quad \mathbf{B} = \begin{bmatrix} 2 \\ 0 \end{bmatrix}, \quad \mathbf{C} = \begin{bmatrix} 1 & 0 \\ 0 & 1 \end{bmatrix}, \quad \mathbf{D} = \begin{bmatrix} 0 \\ 0 \end{bmatrix}
\end{equation}

\textbf{Step 6: Check controllability and observability}
\begin{equation}
\mathcal{C} = \begin{bmatrix} 2 & -4 \\ 0 & 20 \end{bmatrix}, \quad \det(\mathcal{C}) = 40 \neq 0 \quad \Rightarrow \text{Controllable}
\end{equation}
\begin{equation}
\mathcal{O} = \begin{bmatrix} 1 & 0 \\ 0 & 1 \\ -2 & -0.2 \\ 10 & -10 \end{bmatrix}, \quad \text{rank}(\mathcal{O}) = 2 \quad \Rightarrow \text{Observable}
\end{equation}

The DC motor is both controllable and observable---a minimal realization.
\end{example}

\subsection{Example: Transfer Function to Controllable Canonical Form}

\begin{example}[Third-Order System to CCF]
\label{ex:third_order_ccf}
Convert to controllable canonical form:
\begin{equation}
G(s) = \frac{2s^2 + 5s + 3}{s^3 + 4s^2 + 5s + 2}
\end{equation}

\textbf{Step 1: Verify the transfer function is strictly proper}

Numerator degree = 2, denominator degree = 3. Strictly proper. $\checkmark$

\textbf{Step 2: Extract coefficients}

Denominator: $s^3 + 4s^2 + 5s + 2$ $\Rightarrow$ $a_2 = 4$, $a_1 = 5$, $a_0 = 2$

Numerator: $2s^2 + 5s + 3$ $\Rightarrow$ $b_2 = 2$, $b_1 = 5$, $b_0 = 3$

\textbf{Step 3: Write the CCF matrices}
\begin{equation}
\mathbf{A}_c = \begin{bmatrix} 0 & 1 & 0 \\ 0 & 0 & 1 \\ -2 & -5 & -4 \end{bmatrix}, \quad \mathbf{B}_c = \begin{bmatrix} 0 \\ 0 \\ 1 \end{bmatrix}
\end{equation}
\begin{equation}
\mathbf{C}_c = \begin{bmatrix} 3 & 5 & 2 \end{bmatrix}, \quad \mathbf{D}_c = 0
\end{equation}

\textbf{Step 4: Find the poles}
\begin{equation}
s^3 + 4s^2 + 5s + 2 = (s+1)(s^2 + 3s + 2) = (s+1)(s+1)(s+2) = (s+1)^2(s+2)
\end{equation}

Poles: $s = -1$ (repeated), $s = -2$. The system is stable.

\textbf{Step 5: Verify controllability}
\begin{equation}
\mathcal{C} = \begin{bmatrix} 0 & 0 & 1 \\ 0 & 1 & -4 \\ 1 & -4 & 11 \end{bmatrix}
\end{equation}
\begin{equation}
\det(\mathcal{C}) = 0(1 \cdot 11 - (-4)(-4)) - 0 + 1(0 \cdot (-4) - 1 \cdot 0) = -1 \neq 0
\end{equation}

The system is controllable, as guaranteed by CCF.
\end{example}

\subsection{Example: Computing the State Transition Matrix}

\begin{example}[State Transition Matrix for Coupled System]
\label{ex:state_transition}
Compute $e^{\mathbf{A}t}$ for:
\begin{equation}
\mathbf{A} = \begin{bmatrix} -3 & 1 \\ 0 & -2 \end{bmatrix}
\end{equation}

\textbf{Method: Laplace Transform}

\textbf{Step 1: Compute $s\mathbf{I} - \mathbf{A}$}
\begin{equation}
s\mathbf{I} - \mathbf{A} = \begin{bmatrix} s+3 & -1 \\ 0 & s+2 \end{bmatrix}
\end{equation}

\textbf{Step 2: Compute the determinant and adjugate}
\begin{equation}
\det(s\mathbf{I} - \mathbf{A}) = (s+3)(s+2) - 0 = (s+3)(s+2)
\end{equation}
\begin{equation}
\text{adj}(s\mathbf{I} - \mathbf{A}) = \begin{bmatrix} s+2 & 1 \\ 0 & s+3 \end{bmatrix}
\end{equation}

\textbf{Step 3: Compute $(s\mathbf{I} - \mathbf{A})^{-1}$}
\begin{equation}
(s\mathbf{I} - \mathbf{A})^{-1} = \frac{1}{(s+3)(s+2)}\begin{bmatrix} s+2 & 1 \\ 0 & s+3 \end{bmatrix}
\end{equation}

\textbf{Step 4: Partial fractions for each element}

Element (1,1): $\frac{s+2}{(s+3)(s+2)} = \frac{1}{s+3}$

Element (1,2): $\frac{1}{(s+3)(s+2)} = \frac{1}{s+2} - \frac{1}{s+3}$

Element (2,1): $0$

Element (2,2): $\frac{s+3}{(s+3)(s+2)} = \frac{1}{s+2}$

\textbf{Step 5: Inverse Laplace transform}
\begin{equation}
e^{\mathbf{A}t} = \begin{bmatrix} e^{-3t} & e^{-2t} - e^{-3t} \\ 0 & e^{-2t} \end{bmatrix}
\end{equation}

\textbf{Verification:} Check that $e^{\mathbf{A} \cdot 0} = \mathbf{I}$:
\begin{equation}
e^{\mathbf{A} \cdot 0} = \begin{bmatrix} 1 & 1-1 \\ 0 & 1 \end{bmatrix} = \begin{bmatrix} 1 & 0 \\ 0 & 1 \end{bmatrix} = \mathbf{I} \quad \checkmark
\end{equation}
\end{example}

\subsection{Example: Controllability Analysis}

\begin{example}[Controllability of Coupled Tanks]
\label{ex:coupled_tanks_controllability}
Two tanks are connected by a pipe. Tank 1 has cross-sectional area $A_1$ and is fed by pump 1. Tank 2 has area $A_2$ and is fed by pump 2. The connecting pipe has resistance $R_{12}$. Tank 2 drains through resistance $R_2$.

The state-space model (from fluid mass balance) is:
\begin{equation}
\mathbf{A} = \begin{bmatrix} -\frac{1}{A_1 R_{12}} & \frac{1}{A_1 R_{12}} \\ \frac{1}{A_2 R_{12}} & -\frac{1}{A_2 R_{12}} - \frac{1}{A_2 R_2} \end{bmatrix}, \quad \mathbf{B} = \begin{bmatrix} \frac{1}{A_1} & 0 \\ 0 & \frac{1}{A_2} \end{bmatrix}
\end{equation}

where states $x_1, x_2$ are tank levels and inputs $u_1, u_2$ are pump flow rates.

\textbf{Numerical values:} $A_1 = A_2 = 1$\,m$^2$, $R_{12} = 2$\,s/m$^2$, $R_2 = 1$\,s/m$^2$:
\begin{equation}
\mathbf{A} = \begin{bmatrix} -0.5 & 0.5 \\ 0.5 & -1.5 \end{bmatrix}, \quad \mathbf{B} = \begin{bmatrix} 1 & 0 \\ 0 & 1 \end{bmatrix}
\end{equation}

\textbf{Controllability test:}
\begin{equation}
\mathcal{C} = \begin{bmatrix} \mathbf{B} & \mathbf{A}\mathbf{B} \end{bmatrix} = \begin{bmatrix} 1 & 0 & -0.5 & 0.5 \\ 0 & 1 & 0.5 & -1.5 \end{bmatrix}
\end{equation}

Since $\mathbf{B} = \mathbf{I}_{2 \times 2}$ and has rank 2, the controllability matrix has rank 2. The system is \textbf{controllable}.

\textbf{Physical interpretation:} Each pump can directly control its respective tank level. The coupling through $R_{12}$ allows each pump to also influence the other tank indirectly.

\textbf{What if only pump 1 exists?}
\begin{equation}
\mathbf{B}' = \begin{bmatrix} 1 \\ 0 \end{bmatrix}
\end{equation}
\begin{equation}
\mathcal{C}' = \begin{bmatrix} 1 & -0.5 \\ 0 & 0.5 \end{bmatrix}, \quad \det(\mathcal{C}') = 0.5 \neq 0
\end{equation}

Still controllable! The coupling through the connecting pipe allows pump 1 to indirectly control tank 2.
\end{example}

\subsection{Example: Observability Analysis}

\begin{example}[Observability with Different Sensor Configurations]
\label{ex:observability_sensors}
Consider the mass-spring-damper system:
\begin{equation}
\mathbf{A} = \begin{bmatrix} 0 & 1 \\ -4 & -2 \end{bmatrix}, \quad \mathbf{B} = \begin{bmatrix} 0 \\ 1 \end{bmatrix}
\end{equation}

Analyze observability for different sensor configurations.

\textbf{Case 1: Position sensor only}
\begin{equation}
\mathbf{C}_1 = \begin{bmatrix} 1 & 0 \end{bmatrix}
\end{equation}
\begin{equation}
\mathcal{O}_1 = \begin{bmatrix} 1 & 0 \\ 0 & 1 \end{bmatrix}, \quad \det(\mathcal{O}_1) = 1 \neq 0 \quad \Rightarrow \text{Observable}
\end{equation}

\textbf{Case 2: Velocity sensor only}
\begin{equation}
\mathbf{C}_2 = \begin{bmatrix} 0 & 1 \end{bmatrix}
\end{equation}
\begin{equation}
\mathbf{C}_2\mathbf{A} = \begin{bmatrix} 0 & 1 \end{bmatrix}\begin{bmatrix} 0 & 1 \\ -4 & -2 \end{bmatrix} = \begin{bmatrix} -4 & -2 \end{bmatrix}
\end{equation}
\begin{equation}
\mathcal{O}_2 = \begin{bmatrix} 0 & 1 \\ -4 & -2 \end{bmatrix}, \quad \det(\mathcal{O}_2) = 0 - (-4) = 4 \neq 0 \quad \Rightarrow \text{Observable}
\end{equation}

\textbf{Case 3: Acceleration sensor only}

Acceleration is $\ddot{x} = -4x - 2\dot{x} + u = -4x_1 - 2x_2 + u$.

For observability analysis with $u$ known, the effective output is:
\begin{equation}
y = \ddot{x} - u = -4x_1 - 2x_2 \quad \Rightarrow \quad \mathbf{C}_3 = \begin{bmatrix} -4 & -2 \end{bmatrix}
\end{equation}
\begin{equation}
\mathbf{C}_3\mathbf{A} = \begin{bmatrix} -4 & -2 \end{bmatrix}\begin{bmatrix} 0 & 1 \\ -4 & -2 \end{bmatrix} = \begin{bmatrix} 8 & 0 \end{bmatrix}
\end{equation}
\begin{equation}
\mathcal{O}_3 = \begin{bmatrix} -4 & -2 \\ 8 & 0 \end{bmatrix}, \quad \det(\mathcal{O}_3) = 0 - (-16) = 16 \neq 0 \quad \Rightarrow \text{Observable}
\end{equation}

\textbf{Conclusion:} For this coupled system, \textit{any single} measurement (position, velocity, or acceleration) is sufficient for observability because the dynamics couple all states. This is a consequence of the off-diagonal coupling in $\mathbf{A}$.
\end{example}

\subsection{Example: Transformation Between Canonical Forms}

\begin{example}[CCF to OCF Transformation]
\label{ex:ccf_to_ocf}
Transform from controllable canonical form to observable canonical form for:
\begin{equation}
G(s) = \frac{2s + 3}{s^2 + 5s + 6}
\end{equation}

\textbf{CCF:}
\begin{equation}
\mathbf{A}_c = \begin{bmatrix} 0 & 1 \\ -6 & -5 \end{bmatrix}, \quad \mathbf{B}_c = \begin{bmatrix} 0 \\ 1 \end{bmatrix}, \quad \mathbf{C}_c = \begin{bmatrix} 3 & 2 \end{bmatrix}
\end{equation}

\textbf{OCF:}
\begin{equation}
\mathbf{A}_o = \begin{bmatrix} 0 & -6 \\ 1 & -5 \end{bmatrix}, \quad \mathbf{B}_o = \begin{bmatrix} 3 \\ 2 \end{bmatrix}, \quad \mathbf{C}_o = \begin{bmatrix} 0 & 1 \end{bmatrix}
\end{equation}

\textbf{Step 1: Compute controllability matrices}
\begin{equation}
\mathcal{C}_c = \begin{bmatrix} 0 & 1 \\ 1 & -5 \end{bmatrix}, \quad \mathcal{C}_o = \begin{bmatrix} 3 & -12 \\ 2 & -7 \end{bmatrix}
\end{equation}

\textbf{Step 2: Compute $\mathcal{C}_c^{-1}$}
\begin{equation}
\det(\mathcal{C}_c) = -1, \quad \mathcal{C}_c^{-1} = \begin{bmatrix} 5 & 1 \\ 1 & 0 \end{bmatrix}
\end{equation}

\textbf{Step 3: Find the transformation matrix}
\begin{equation}
\mathbf{T} = \mathcal{C}_o \mathcal{C}_c^{-1} = \begin{bmatrix} 3 & -12 \\ 2 & -7 \end{bmatrix}\begin{bmatrix} 5 & 1 \\ 1 & 0 \end{bmatrix} = \begin{bmatrix} 3 & 3 \\ 3 & 2 \end{bmatrix}
\end{equation}

\textbf{Step 4: Verify}
\begin{equation}
\mathbf{T}\mathbf{A}_c\mathbf{T}^{-1} = \begin{bmatrix} 3 & 3 \\ 3 & 2 \end{bmatrix}\begin{bmatrix} 0 & 1 \\ -6 & -5 \end{bmatrix}\begin{bmatrix} -\frac{2}{3} & 1 \\ 1 & -1 \end{bmatrix}
\end{equation}

Let me compute $\mathbf{T}^{-1}$:
\begin{equation}
\det(\mathbf{T}) = 6 - 9 = -3, \quad \mathbf{T}^{-1} = -\frac{1}{3}\begin{bmatrix} 2 & -3 \\ -3 & 3 \end{bmatrix} = \begin{bmatrix} -\frac{2}{3} & 1 \\ 1 & -1 \end{bmatrix}
\end{equation}

Continuing:
\begin{equation}
\mathbf{T}\mathbf{A}_c = \begin{bmatrix} 3 & 3 \\ 3 & 2 \end{bmatrix}\begin{bmatrix} 0 & 1 \\ -6 & -5 \end{bmatrix} = \begin{bmatrix} -18 & -12 \\ -12 & -7 \end{bmatrix}
\end{equation}
\begin{equation}
\mathbf{T}\mathbf{A}_c\mathbf{T}^{-1} = \begin{bmatrix} -18 & -12 \\ -12 & -7 \end{bmatrix}\begin{bmatrix} -\frac{2}{3} & 1 \\ 1 & -1 \end{bmatrix} = \begin{bmatrix} 0 & -6 \\ 1 & -5 \end{bmatrix} = \mathbf{A}_o \quad \checkmark
\end{equation}
\end{example}

\subsection{Example: MIMO State-Space Representation}

\begin{example}[Two-Input Two-Output Thermal System]
\label{ex:mimo_thermal}
A two-zone building has temperatures $T_1$ and $T_2$, with heaters providing heat inputs $Q_1$ and $Q_2$. The zones are thermally coupled.

\textbf{Energy balance:}
\begin{align}
C_1\dot{T}_1 &= Q_1 - \frac{T_1 - T_{\text{amb}}}{R_1} - \frac{T_1 - T_2}{R_{12}} \\
C_2\dot{T}_2 &= Q_2 - \frac{T_2 - T_{\text{amb}}}{R_2} - \frac{T_2 - T_1}{R_{12}}
\end{align}

Using deviation variables $x_i = T_i - T_{\text{amb}}$ and simplifying:
\begin{equation}
\mathbf{A} = \begin{bmatrix} -\frac{1}{C_1}\left(\frac{1}{R_1} + \frac{1}{R_{12}}\right) & \frac{1}{C_1 R_{12}} \\ \frac{1}{C_2 R_{12}} & -\frac{1}{C_2}\left(\frac{1}{R_2} + \frac{1}{R_{12}}\right) \end{bmatrix}
\end{equation}
\begin{equation}
\mathbf{B} = \begin{bmatrix} \frac{1}{C_1} & 0 \\ 0 & \frac{1}{C_2} \end{bmatrix}, \quad \mathbf{C} = \begin{bmatrix} 1 & 0 \\ 0 & 1 \end{bmatrix}, \quad \mathbf{D} = \begin{bmatrix} 0 & 0 \\ 0 & 0 \end{bmatrix}
\end{equation}

\textbf{Numerical values:} $C_1 = C_2 = 1000$\,J/K, $R_1 = R_2 = 0.01$\,K/W, $R_{12} = 0.05$\,K/W:
\begin{equation}
\mathbf{A} = \begin{bmatrix} -0.12 & 0.02 \\ 0.02 & -0.12 \end{bmatrix}, \quad \mathbf{B} = \begin{bmatrix} 0.001 & 0 \\ 0 & 0.001 \end{bmatrix}
\end{equation}

\textbf{Transfer function matrix:}
\begin{equation}
\mathbf{G}(s) = \mathbf{C}(s\mathbf{I} - \mathbf{A})^{-1}\mathbf{B}
\end{equation}

Let me compute $(s\mathbf{I} - \mathbf{A})^{-1}$:
\begin{equation}
s\mathbf{I} - \mathbf{A} = \begin{bmatrix} s + 0.12 & -0.02 \\ -0.02 & s + 0.12 \end{bmatrix}
\end{equation}
\begin{equation}
\det(s\mathbf{I} - \mathbf{A}) = (s + 0.12)^2 - 0.0004 = s^2 + 0.24s + 0.014 = (s + 0.1)(s + 0.14)
\end{equation}

The poles are $s = -0.1$ and $s = -0.14$, corresponding to time constants of 10 s and 7.14 s.

\textbf{Interpretation of the transfer function matrix:}
\begin{equation}
\mathbf{G}(s) = \begin{bmatrix} G_{11}(s) & G_{12}(s) \\ G_{21}(s) & G_{22}(s) \end{bmatrix}
\end{equation}

\begin{itemize}
    \item $G_{11}(s)$, $G_{22}(s)$: Direct channels (heater to its own zone)
    \item $G_{12}(s)$, $G_{21}(s)$: Cross-coupling (heater affecting the other zone)
\end{itemize}

Due to symmetry, $G_{11} = G_{22}$ and $G_{12} = G_{21}$.
\end{example}

\subsection{Example: Modal Analysis Using State-Space}

\begin{example}[Modal Decomposition of a Vibrating System]
\label{ex:modal_analysis}
A two-degree-of-freedom mechanical system has:
\begin{equation}
\mathbf{A} = \begin{bmatrix} 0 & 0 & 1 & 0 \\ 0 & 0 & 0 & 1 \\ -3 & 1 & -0.4 & 0.1 \\ 1 & -2 & 0.1 & -0.3 \end{bmatrix}
\end{equation}

This represents two coupled masses with states $[x_1, x_2, \dot{x}_1, \dot{x}_2]^\top$.

\textbf{Step 1: Find eigenvalues}

The characteristic polynomial is:
\begin{equation}
\det(s\mathbf{I} - \mathbf{A}) = s^4 + 0.7s^3 + 4.97s^2 + 1.4s + 5 = 0
\end{equation}

Eigenvalues (computed numerically):
\begin{align}
\lambda_{1,2} &= -0.175 \pm 1.38j \quad \text{(Mode 1: lower frequency)} \\
\lambda_{3,4} &= -0.175 \pm 2.05j \quad \text{(Mode 2: higher frequency)}
\end{align}

\textbf{Step 2: Interpret the modes}

\textit{Mode 1} ($\omega_1 \approx 1.38$ rad/s): Lower frequency oscillation. From eigenvector analysis, both masses move in phase (in the same direction).

\textit{Mode 2} ($\omega_2 \approx 2.05$ rad/s): Higher frequency oscillation. The masses move in opposite directions (out of phase).

\textbf{Step 3: Damping ratios}

For both modes, the real part is $-0.175$ and the imaginary parts are $\pm\omega_d$.

Mode 1: $\zeta_1 = \frac{0.175}{\sqrt{0.175^2 + 1.38^2}} \approx 0.126$

Mode 2: $\zeta_2 = \frac{0.175}{\sqrt{0.175^2 + 2.05^2}} \approx 0.085$

Both modes are lightly damped, with the higher-frequency mode being less damped.

\textbf{Modal state-space form:}

By transforming to modal coordinates using the eigenvector matrix, we can decouple the system into independent second-order oscillators, making analysis and control design more intuitive.
\end{example}

%=============================================================================
\section{Algorithm Implementations}
\label{sec:algorithms}
%=============================================================================

This section provides production-ready algorithms for the key computations covered in this chapter. Each algorithm includes complexity analysis and implementation guidance.

\subsection{Transfer Function to Controllable Canonical Form}

\begin{algorithm}[H]
\caption{Transfer Function to Controllable Canonical Form}
\label{alg:tf_to_ccf}
\begin{algorithmic}[1]
\REQUIRE Numerator coefficients $[b_0, b_1, \ldots, b_{n-1}]$, denominator coefficients $[a_0, a_1, \ldots, a_{n-1}]$ (monic: $a_n = 1$)
\ENSURE State-space matrices $\mathbf{A}, \mathbf{B}, \mathbf{C}, \mathbf{D}$
\STATE $n \gets$ length of denominator $- 1$ \COMMENT{System order}
\STATE Initialize $\mathbf{A}$ as $n \times n$ zero matrix
\STATE Initialize $\mathbf{B}$ as $n \times 1$ zero vector
\STATE
\COMMENT{Build companion matrix structure}
\FOR{$i = 1$ \TO $n-1$}
    \STATE $\mathbf{A}[i, i+1] \gets 1$ \COMMENT{Superdiagonal ones}
\ENDFOR
\FOR{$j = 1$ \TO $n$}
    \STATE $\mathbf{A}[n, j] \gets -a_{j-1}$ \COMMENT{Last row: negated denominator coefficients}
\ENDFOR
\STATE
\STATE $\mathbf{B}[n] \gets 1$ \COMMENT{Input enters last state}
\STATE $\mathbf{C} \gets [b_0, b_1, \ldots, b_{n-1}]$ \COMMENT{Numerator coefficients}
\STATE $\mathbf{D} \gets 0$
\STATE
\RETURN $\mathbf{A}, \mathbf{B}, \mathbf{C}, \mathbf{D}$
\end{algorithmic}
\end{algorithm}

\textbf{Computational complexity:} $\mathcal{O}(n)$ operations.

\textbf{Implementation notes:}
\begin{itemize}
    \item Ensure the denominator is monic (leading coefficient = 1)
    \item For proper (not strictly proper) transfer functions, first perform polynomial long division to extract $\mathbf{D}$
    \item Verify by computing the transfer function from the resulting state-space model
\end{itemize}

\subsection{Controllability Matrix and Test}

\begin{algorithm}[H]
\caption{Controllability Matrix Computation and Test}
\label{alg:controllability}
\begin{algorithmic}[1]
\REQUIRE State matrix $\mathbf{A} \in \mathbb{R}^{n \times n}$, input matrix $\mathbf{B} \in \mathbb{R}^{n \times m}$
\ENSURE Controllability matrix $\mathcal{C}$, boolean controllability flag
\STATE $n \gets$ number of rows in $\mathbf{A}$
\STATE $m \gets$ number of columns in $\mathbf{B}$
\STATE Initialize $\mathcal{C}$ as $n \times (nm)$ matrix
\STATE
\STATE $\mathcal{C}[:, 1:m] \gets \mathbf{B}$ \COMMENT{First block is $\mathbf{B}$}
\STATE $\mathbf{A}^k\mathbf{B} \gets \mathbf{B}$
\FOR{$k = 1$ \TO $n-1$}
    \STATE $\mathbf{A}^k\mathbf{B} \gets \mathbf{A} \cdot \mathbf{A}^{k-1}\mathbf{B}$ \COMMENT{Compute $\mathbf{A}^k\mathbf{B}$}
    \STATE $\mathcal{C}[:, km+1:(k+1)m] \gets \mathbf{A}^k\mathbf{B}$
\ENDFOR
\STATE
\STATE $r \gets \text{rank}(\mathcal{C})$ \COMMENT{Numerical rank computation}
\STATE \textbf{controllable} $\gets (r = n)$
\STATE
\RETURN $\mathcal{C}$, \textbf{controllable}
\end{algorithmic}
\end{algorithm}

\textbf{Computational complexity:} $\mathcal{O}(n^3)$ for the matrix multiplications and rank computation.

\textbf{Numerical considerations:}
\begin{itemize}
    \item Use SVD for robust rank computation: count singular values above a threshold
    \item Typical threshold: $\max(\sigma_i) \cdot n \cdot \epsilon_{\text{machine}}$
    \item For large systems, the controllability Gramian approach is more numerically stable
\end{itemize}

\subsection{State Transition Matrix via Series Expansion}

\begin{algorithm}[H]
\caption{Matrix Exponential via Taylor Series}
\label{alg:matrix_exp_series}
\begin{algorithmic}[1]
\REQUIRE State matrix $\mathbf{A} \in \mathbb{R}^{n \times n}$, time $t$, tolerance $\epsilon$
\ENSURE State transition matrix $e^{\mathbf{A}t}$
\STATE $\mathbf{\Phi} \gets \mathbf{I}_n$ \COMMENT{Initialize to identity}
\STATE $\mathbf{Term} \gets \mathbf{I}_n$ \COMMENT{Current term in series}
\STATE $k \gets 0$
\REPEAT
    \STATE $k \gets k + 1$
    \STATE $\mathbf{Term} \gets \mathbf{Term} \cdot \mathbf{A}t / k$ \COMMENT{$(\mathbf{A}t)^k / k!$}
    \STATE $\mathbf{\Phi} \gets \mathbf{\Phi} + \mathbf{Term}$
\UNTIL{$\|\mathbf{Term}\|_F < \epsilon$}
\STATE
\RETURN $\mathbf{\Phi}$
\end{algorithmic}
\end{algorithm}

\textbf{Warning:} This naive implementation can suffer from numerical issues for large $\|At\|$. Use scaling and squaring methods (as in MATLAB's \texttt{expm}) for production code.

\subsection{State Transition Matrix via Eigendecomposition}

\begin{algorithm}[H]
\caption{Matrix Exponential via Eigendecomposition}
\label{alg:matrix_exp_eigen}
\begin{algorithmic}[1]
\REQUIRE State matrix $\mathbf{A} \in \mathbb{R}^{n \times n}$, time $t$
\ENSURE State transition matrix $e^{\mathbf{A}t}$
\STATE Compute eigendecomposition: $\mathbf{A} = \mathbf{V}\boldsymbol{\Lambda}\mathbf{V}^{-1}$
\STATE \COMMENT{$\boldsymbol{\Lambda} = \text{diag}(\lambda_1, \ldots, \lambda_n)$, $\mathbf{V}$ has eigenvectors as columns}
\IF{$\mathbf{A}$ is not diagonalizable}
    \STATE Use Jordan decomposition or scaling-squaring method instead
\ENDIF
\STATE
\STATE Compute diagonal exponential: $e^{\boldsymbol{\Lambda}t} = \text{diag}(e^{\lambda_1 t}, \ldots, e^{\lambda_n t})$
\STATE $\mathbf{\Phi} \gets \mathbf{V} \cdot e^{\boldsymbol{\Lambda}t} \cdot \mathbf{V}^{-1}$
\STATE
\RETURN $\mathbf{\Phi}$
\end{algorithmic}
\end{algorithm}

\textbf{Computational complexity:} $\mathcal{O}(n^3)$ for eigendecomposition.

\textbf{Limitations:}
\begin{itemize}
    \item Requires distinct eigenvalues (or use Jordan form for repeated eigenvalues)
    \item Eigenvector computation can be numerically sensitive
    \item For complex eigenvalues, the result will be real but computation involves complex arithmetic
\end{itemize}

\subsection{Similarity Transformation}

\begin{algorithm}[H]
\caption{Apply Similarity Transformation}
\label{alg:similarity}
\begin{algorithmic}[1]
\REQUIRE Original system $(\mathbf{A}, \mathbf{B}, \mathbf{C}, \mathbf{D})$, invertible transformation $\mathbf{T}$
\ENSURE Transformed system $(\bar{\mathbf{A}}, \bar{\mathbf{B}}, \bar{\mathbf{C}}, \bar{\mathbf{D}})$
\STATE Verify $\det(\mathbf{T}) \neq 0$ \COMMENT{Check invertibility}
\STATE Compute $\mathbf{T}^{-1}$
\STATE
\STATE $\bar{\mathbf{A}} \gets \mathbf{T} \cdot \mathbf{A} \cdot \mathbf{T}^{-1}$
\STATE $\bar{\mathbf{B}} \gets \mathbf{T} \cdot \mathbf{B}$
\STATE $\bar{\mathbf{C}} \gets \mathbf{C} \cdot \mathbf{T}^{-1}$
\STATE $\bar{\mathbf{D}} \gets \mathbf{D}$ \COMMENT{Unchanged}
\STATE
\RETURN $(\bar{\mathbf{A}}, \bar{\mathbf{B}}, \bar{\mathbf{C}}, \bar{\mathbf{D}})$
\end{algorithmic}
\end{algorithm}

\textbf{Verification:} After transformation, verify that eigenvalues of $\bar{\mathbf{A}}$ match those of $\mathbf{A}$.

\subsection{State-Space to Transfer Function}

\begin{algorithm}[H]
\caption{State-Space to Transfer Function (SISO)}
\label{alg:ss_to_tf}
\begin{algorithmic}[1]
\REQUIRE State-space matrices $\mathbf{A} \in \mathbb{R}^{n \times n}$, $\mathbf{B} \in \mathbb{R}^{n \times 1}$, $\mathbf{C} \in \mathbb{R}^{1 \times n}$, $\mathbf{D}$ (scalar)
\ENSURE Numerator polynomial $\mathbf{num}$, denominator polynomial $\mathbf{den}$
\STATE
\COMMENT{Compute characteristic polynomial (denominator)}
\STATE $\mathbf{den} \gets \text{poly}(\mathbf{A})$ \COMMENT{Coefficients of $\det(s\mathbf{I} - \mathbf{A})$}
\STATE
\COMMENT{Compute numerator via adjugate formula}
\STATE Form $(s\mathbf{I} - \mathbf{A})$ symbolically
\STATE Compute $\text{adj}(s\mathbf{I} - \mathbf{A})$ \COMMENT{Cofactor matrix transpose}
\STATE $\mathbf{num}(s) \gets \mathbf{C} \cdot \text{adj}(s\mathbf{I} - \mathbf{A}) \cdot \mathbf{B} + \mathbf{D} \cdot \det(s\mathbf{I} - \mathbf{A})$
\STATE Extract polynomial coefficients of $\mathbf{num}(s)$
\STATE
\RETURN $\mathbf{num}$, $\mathbf{den}$
\end{algorithmic}
\end{algorithm}

\textbf{Practical note:} For numerical computation, use the formula:
\begin{equation}
G(s) = \frac{\det\begin{bmatrix} s\mathbf{I} - \mathbf{A} & \mathbf{B} \\ -\mathbf{C} & \mathbf{D} \end{bmatrix}}{\det(s\mathbf{I} - \mathbf{A})}
\end{equation}

%=============================================================================
\section{Chapter Summary}
\label{sec:chapter_summary}
%=============================================================================

This chapter introduced state-space representation, the foundation for modern control theory. Here are the key concepts and their significance.

\subsection{Key Concepts}

\textbf{State-Space Fundamentals:}
\begin{itemize}
    \item The state vector $\mathbf{x}(t)$ captures all information needed to predict future system behavior
    \item State-space form: $\dot{\mathbf{x}} = \mathbf{A}\mathbf{x} + \mathbf{B}\mathbf{u}$, $\mathbf{y} = \mathbf{C}\mathbf{x} + \mathbf{D}\mathbf{u}$
    \item The $\mathbf{A}$ matrix determines internal dynamics; its eigenvalues are the system poles
    \item Transfer function: $G(s) = \mathbf{C}(s\mathbf{I} - \mathbf{A})^{-1}\mathbf{B} + \mathbf{D}$
\end{itemize}

\textbf{Canonical Forms:}
\begin{itemize}
    \item \textit{Controllable canonical form:} Characteristic polynomial in last row of $\mathbf{A}$; guarantees controllability
    \item \textit{Observable canonical form:} Characteristic polynomial in last column of $\mathbf{A}$; guarantees observability
    \item \textit{Diagonal form:} Decoupled modes with eigenvalues on diagonal; requires distinct poles
    \item \textit{Jordan form:} Handles repeated eigenvalues with Jordan blocks
\end{itemize}

\textbf{State Transition Matrix:}
\begin{itemize}
    \item $\boldsymbol{\Phi}(t) = e^{\mathbf{A}t}$ describes free response evolution
    \item Complete solution: $\mathbf{x}(t) = e^{\mathbf{A}t}\mathbf{x}_0 + \int_0^t e^{\mathbf{A}(t-\tau)}\mathbf{B}\mathbf{u}(\tau)\,d\tau$
    \item Computation methods: Laplace transform, eigendecomposition, series expansion
\end{itemize}

\textbf{Controllability:}
\begin{itemize}
    \item A system is controllable if any state can be reached from any initial state using inputs
    \item Test: rank$(\mathcal{C}) = n$ where $\mathcal{C} = [\mathbf{B} \; \mathbf{AB} \; \cdots \; \mathbf{A}^{n-1}\mathbf{B}]$
    \item Physical meaning: Can the input influence all system modes?
    \item The controllability Gramian $\mathbf{W}_c$ provides quantitative controllability measures
\end{itemize}

\textbf{Observability:}
\begin{itemize}
    \item A system is observable if the initial state can be determined from output measurements
    \item Test: rank$(\mathcal{O}) = n$ where $\mathcal{O} = [\mathbf{C}^\top \; \mathbf{A}^\top\mathbf{C}^\top \; \cdots \; (\mathbf{A}^\top)^{n-1}\mathbf{C}^\top]^\top$
    \item Physical meaning: Do all system modes affect the measurements?
    \item Duality: Observability of $(\mathbf{A}, \mathbf{C})$ $\Leftrightarrow$ controllability of $(\mathbf{A}^\top, \mathbf{C}^\top)$
\end{itemize}

\textbf{Similarity Transformations:}
\begin{itemize}
    \item New coordinates: $\bar{\mathbf{x}} = \mathbf{T}\mathbf{x}$ yields $\bar{\mathbf{A}} = \mathbf{T}\mathbf{A}\mathbf{T}^{-1}$, $\bar{\mathbf{B}} = \mathbf{T}\mathbf{B}$, $\bar{\mathbf{C}} = \mathbf{C}\mathbf{T}^{-1}$
    \item Preserves: transfer function, eigenvalues, controllability, observability
    \item Used to convert between canonical forms
\end{itemize}

\textbf{Minimal Realization:}
\begin{itemize}
    \item A realization is minimal if it is both controllable and observable
    \item Minimal realizations have the fewest states possible for a given transfer function
    \item Non-minimal realizations have hidden modes (pole-zero cancellations in transfer function)
\end{itemize}

\subsection{Design Guidelines}

\begin{enumerate}
    \item \textbf{Choose canonical form based on application:}
    \begin{itemize}
        \item State feedback design $\rightarrow$ controllable canonical form
        \item Observer design $\rightarrow$ observable canonical form
        \item Modal analysis $\rightarrow$ diagonal form
    \end{itemize}

    \item \textbf{Always verify controllability and observability:}
    \begin{itemize}
        \item Uncontrollable modes cannot be stabilized by feedback
        \item Unobservable modes cannot be estimated from measurements
        \item Both must hold for observer-based control
    \end{itemize}

    \item \textbf{Use physical state variables when possible:}
    \begin{itemize}
        \item Provides physical intuition
        \item Facilitates debugging and validation
        \item Convert to canonical form only when needed for design
    \end{itemize}

    \item \textbf{For numerical computation:}
    \begin{itemize}
        \item Use Gramians rather than rank tests for large systems
        \item Employ balanced realizations for numerical stability
        \item Verify transformations by checking eigenvalue preservation
    \end{itemize}
\end{enumerate}

\subsection{Preview of Coming Topics}

State-space representation enables the modern control methods we will study in subsequent chapters:

\begin{itemize}
    \item \textbf{Chapter 11 (Observer Design):} Design state estimators that reconstruct unmeasured states from output measurements---requires observability

    \item \textbf{Chapter 12 (MIMO Control):} State-space naturally handles multiple inputs and outputs, enabling decoupling and multivariable controller design

    \item \textbf{Chapter 14 (Optimal Control):} Linear Quadratic Regulator (LQR) minimizes a cost function using full state feedback---requires controllability

    \item \textbf{Part V (ML-Based Control):} Machine learning methods for control operate on state representations, building on the foundations established here
\end{itemize}

\subsection{Key Equations Summary}

\begin{center}
\renewcommand{\arraystretch}{1.5}
\begin{tabular}{ll}
\toprule
\textbf{Concept} & \textbf{Equation} \\
\midrule
State-space model & $\dot{\mathbf{x}} = \mathbf{A}\mathbf{x} + \mathbf{B}\mathbf{u}$, \; $\mathbf{y} = \mathbf{C}\mathbf{x} + \mathbf{D}\mathbf{u}$ \\
Transfer function & $G(s) = \mathbf{C}(s\mathbf{I} - \mathbf{A})^{-1}\mathbf{B} + \mathbf{D}$ \\
State transition matrix & $\boldsymbol{\Phi}(t) = e^{\mathbf{A}t} = \mathcal{L}^{-1}\{(s\mathbf{I} - \mathbf{A})^{-1}\}$ \\
Complete solution & $\mathbf{x}(t) = e^{\mathbf{A}t}\mathbf{x}_0 + \int_0^t e^{\mathbf{A}(t-\tau)}\mathbf{B}\mathbf{u}(\tau)\,d\tau$ \\
Controllability matrix & $\mathcal{C} = [\mathbf{B} \; \mathbf{AB} \; \mathbf{A}^2\mathbf{B} \; \cdots \; \mathbf{A}^{n-1}\mathbf{B}]$ \\
Controllability test & System controllable $\Leftrightarrow$ rank$(\mathcal{C}) = n$ \\
Observability matrix & $\mathcal{O} = [\mathbf{C}^\top \; \mathbf{A}^\top\mathbf{C}^\top \; \cdots \; (\mathbf{A}^\top)^{n-1}\mathbf{C}^\top]^\top$ \\
Observability test & System observable $\Leftrightarrow$ rank$(\mathcal{O}) = n$ \\
Similarity transformation & $\bar{\mathbf{A}} = \mathbf{T}\mathbf{A}\mathbf{T}^{-1}$, $\bar{\mathbf{B}} = \mathbf{T}\mathbf{B}$, $\bar{\mathbf{C}} = \mathbf{C}\mathbf{T}^{-1}$ \\
Controllability Gramian & $\mathbf{A}\mathbf{W}_c + \mathbf{W}_c\mathbf{A}^\top + \mathbf{B}\mathbf{B}^\top = \mathbf{0}$ \\
Observability Gramian & $\mathbf{A}^\top\mathbf{W}_o + \mathbf{W}_o\mathbf{A} + \mathbf{C}^\top\mathbf{C} = \mathbf{0}$ \\
\bottomrule
\end{tabular}
\end{center}

